% ============================================================
% RAPPORT PFE — LogiWay
% Plateforme Intelligente de Gestion de Flotte
% Version 2 — Organisé en 6 Sprints Agile Scrum
% Compile avec : pdflatex rapport_v2.tex (x2)
% ============================================================

\documentclass[12pt,a4paper]{report}

\usepackage[utf8]{inputenc}
\usepackage[T1]{fontenc}
\usepackage{lmodern}
\usepackage{textcomp}
\usepackage[french]{babel}
\usepackage{pdflscape}
\usepackage{geometry}
\geometry{top=2.5cm,bottom=2.5cm,left=3cm,right=2.5cm,headheight=80pt}
\usepackage{setspace}
\setstretch{1.15}
\setlength{\parskip}{0.35em}
\setlength{\parindent}{0pt}
\setlength{\emergencystretch}{3em}
\tolerance=1200
\hbadness=2000
\usepackage{graphicx}
\usepackage[table,xcdraw]{xcolor}
\usepackage{tikz}
\usepackage{pgfplots}
\pgfplotsset{compat=1.18}
\usepackage{tabularx}
\usepackage{longtable}
\usepackage{booktabs}
\usepackage{multirow}
\usepackage{array}
\usepackage{float}
\usepackage{caption}
\usepackage{subcaption}
\usepackage{enumitem}
\setlist{itemsep=0.15em,topsep=0.2em}
\usepackage{amsmath}
\usepackage{amssymb}
\usepackage{microtype}
\DisableLigatures{encoding = *, family = *}
\usepackage{listings}
\usepackage{tcolorbox}
\tcbuselibrary{listings,breakable,skins}
\usepackage{fancyhdr}
\usepackage{titlesec}
\usepackage{tocloft}
\usepackage{etoolbox}
\usepackage{hyperref}

% ============================================================
% COULEURS
% ============================================================
\definecolor{rougeEsprit}{RGB}{178,0,30}
\definecolor{primary}{RGB}{178,0,30}
\definecolor{noir}{RGB}{0,0,0}
\definecolor{titleColor}{RGB}{0,0,0}
\definecolor{sectionColor}{RGB}{0,0,0}
\definecolor{subsectionColor}{RGB}{0,0,0}
\definecolor{amber}{RGB}{217,119,6}
\definecolor{tablehead}{RGB}{11,31,58}
\definecolor{tablerow}{RGB}{241,245,249}
\definecolor{codebg}{RGB}{30,41,59}
\definecolor{codetext}{RGB}{226,232,240}
\definecolor{lightgray}{RGB}{248,250,252}
\definecolor{navy}{RGB}{11,31,58}
\definecolor{teal}{RGB}{0,128,128}

% ============================================================
% EN-TÊTES / PIEDS
% ============================================================
\pagestyle{fancy}
\fancyhf{}
\fancyhead[L]{\includegraphics[height=1.4cm]{logo_esprit.jpg}}
\fancyhead[R]{\textcolor{rougeEsprit}{\small\nouppercase{\leftmark}}}
\fancyfoot[C]{\textcolor{rougeEsprit}{\thepage}}
\renewcommand{\headrulewidth}{0.5pt}
\renewcommand{\headrule}{%
  \hbox to\headwidth{\color{rougeEsprit}\leaders\hrule height\headrulewidth\hfill}}
\renewcommand{\footrulewidth}{0.4pt}
\renewcommand{\footrule}{%
  \hbox to\headwidth{\color{rougeEsprit}\leaders\hrule height\footrulewidth\hfill}}

\fancypagestyle{plain}{%
  \fancyhf{}
  \fancyfoot[C]{\textcolor{rougeEsprit}{\thepage}}
  \renewcommand{\headrulewidth}{0pt}
  \renewcommand{\footrulewidth}{0pt}
}

% ============================================================
% TITRES
% ============================================================
\titleformat{\chapter}[display]
  {\normalfont\huge\bfseries\color{noir}}
  {\chaptertitlename~\thechapter}{18pt}{\Huge}
\titlespacing*{\chapter}{0pt}{-20pt}{30pt}

\titleformat{\section}
  {\normalfont\Large\bfseries\color{noir}}{\thesection}{1em}{}
\titleformat{\subsection}
  {\normalfont\large\bfseries\color{noir}}{\thesubsection}{1em}{}
\titleformat{\subsubsection}
  {\normalfont\normalsize\bfseries\color{noir}}{\thesubsubsection}{1em}{}

% ============================================================
% TABLEAUX
% ============================================================
\newcolumntype{Y}{>{\raggedright\arraybackslash}X}
\newcolumntype{C}{>{\centering\arraybackslash}X}
\captionsetup{labelfont=bf,textfont=it,skip=6pt}
\captionsetup[table]{position=top}
\captionsetup[figure]{position=bottom}
\newcommand{\thead}[1]{\cellcolor{tablehead}\textcolor{white}{\textbf{#1}}}

% ============================================================
% BOITES
% ============================================================
\tcbset{
  ruleframe/.style={
    enhanced,breakable,colback=lightgray,colframe=rougeEsprit,
    boxrule=0.8pt,arc=4pt,left=6pt,right=6pt,top=4pt,bottom=4pt
  },
  alertbox/.style={
    enhanced,colback=amber!10,colframe=amber,boxrule=1pt,arc=4pt,
    left=6pt,right=6pt,top=4pt,bottom=4pt
  },
  keybox/.style={
    enhanced,colback=navy!5,colframe=navy,boxrule=0.8pt,arc=4pt,
    left=6pt,right=6pt,top=4pt,bottom=4pt
  }
}

% ============================================================
% CODE
% ============================================================
\lstdefinestyle{logicode}{
  backgroundcolor=\color{codebg},
  basicstyle=\ttfamily\small\color{codetext},
  keywordstyle=\color{rougeEsprit}\bfseries,
  commentstyle=\color{gray}\itshape,
  stringstyle=\color{amber},
  numbers=left,numberstyle=\tiny\color{gray},
  stepnumber=1,numbersep=8pt,
  breaklines=true,breakatwhitespace=true,
  frame=none,rulecolor=\color{rougeEsprit},
  showstringspaces=false,tabsize=2,
  captionpos=b
}
\lstset{style=logicode}

% ============================================================
% HYPERLIENS
% ============================================================
\hypersetup{
  colorlinks=true,linkcolor=noir,urlcolor=rougeEsprit,
  citecolor=rougeEsprit,pdftitle={Rapport PFE LogiWay},
  pdfauthor={Kossay Boubaker},
  pdfkeywords={Spring Boot,Angular,Keycloak,GPS,IA,ML,Gestion Flotte}
}

% ============================================================
% COMMANDES PERSONNALISÉES
% ============================================================
\newcommand{\kpi}[2]{\textcolor{rougeEsprit}{\textbf{#1}} : #2}
\newcounter{planitemcounter}

\newcommand{\planitem}[2]{%
  \stepcounter{planitemcounter}%
  \noindent
  \textcolor{navy}{\textbf{\theplanitemcounter}}\hspace{0.4cm}%
  \textcolor{noir}{\textbf{#1}}%
  \leavevmode\leaders\hbox to 1em{\hss.\hss}\hfill
  \textcolor{navy}{\textbf{\pageref{#2}}}\\[0.25cm]%
}

\newcommand{\bandeauchapitre}[1]{%
  \thispagestyle{fancy}%
  \begin{center}
    \rule{3.2cm}{2.6pt}\quad{\itshape\large Chapitre \thechapter}\quad\rule{3.2cm}{2.6pt}\\[0.35cm]
    \textcolor{noir}{\rule{0.9\linewidth}{0.6pt}}\\[0.45cm]
    {\fontsize{22}{26}\selectfont\bfseries\MakeUppercase{#1}}\\[0.35cm]
    \textcolor{noir}{\rule{0.9\linewidth}{0.6pt}}
  \end{center}
  \vspace{1cm}
}

\newenvironment{plandechapitre}{%
  \setcounter{planitemcounter}{0}%
  \noindent{\Large\bfseries Plan}\par\vspace{0.4cm}
}{%
  \vspace{0.6cm}
  \newpage
}

\newcommand{\chapitreavecplan}[1]{%
  \stepcounter{chapter}%
  \markboth{\MakeUppercase{#1}}{\MakeUppercase{#1}}%
  \phantomsection
  \addcontentsline{toc}{chapter}{\protect\numberline{\thechapter}#1}%
  \pdfbookmark[0]{\thechapter\ #1}{ch:\thechapter}%
  \bandeauchapitre{#1}%
}

% ============================================================
% DOCUMENT
% ============================================================
\begin{document}

\pagenumbering{roman}
\pagestyle{plain}

% ────────────────────────────────────────────────────────────
% DÉDICACES
% ────────────────────────────────────────────────────────────
\chapter*{Dédicaces}
\addcontentsline{toc}{chapter}{Dédicaces}

\vspace{0.8cm}

\begin{center}
\begin{minipage}{0.85\textwidth}
\centering
\itshape

À ma chère mère \textbf{Monia},\\[2pt]
pour tous les sacrifices qu'elle a faits pour moi, toute la confiance\\
qu'elle me porte et tout l'amour dont elle m'entoure.

\vspace{0.9cm}

À mon cher père \textbf{Lotfi},\\[2pt]
pour ses encouragements, sa confiance et son soutien moral,\\
et pour son amour infini.

\vspace{0.9cm}

Une pensée toute particulière et une sincère reconnaissance à\\
\textbf{Monsieur Hadhri Hazem}, Chef de Projet chez ExypnoTech Engineering
Services,\\
pour son soutien indéfectible, sa disponibilité de chaque instant et sa
motivation communicative\\
qui m'ont accompagné et porté tout au long de la réalisation de ce projet.

\vspace{0.9cm}

Je dédie également ce travail, avec toute ma gratitude, à tout esprit libre\\
qui a travaillé dur et a cru en lui-même.\\
Que Dieu vous protège et vous garde.

\end{minipage}
\end{center}

\vspace{1.1cm}
{\color{rougeEsprit}\hrule height 0.6pt}
\vspace{1cm}

\begin{center}
\begin{minipage}{0.85\textwidth}
\centering
\itshape

To my supportive half, you have always been on my side.\\[8pt]
To my late friend, you remain within me — I wish time could bring you back.

\end{minipage}
\end{center}

% ────────────────────────────────────────────────────────────
% REMERCIEMENTS
% ────────────────────────────────────────────────────────────
\chapter*{Remerciements}
\addcontentsline{toc}{chapter}{Remerciements}

Je tiens à exprimer ma profonde gratitude à mon encadrant professionnel
\textbf{M.~Fehmi Kharroubi} et au chef de projet \textbf{M.~Hazem Hadhri}
pour leur supervision, leurs conseils précieux et leur soutien constant tout
au long de ce projet.

Mes remerciements vont également à toute l'équipe d'\textbf{ExypnoTech
Engineering Services} pour l'accueil chaleureux et l'accompagnement technique.

Je remercie sincèrement mes enseignants d'\textbf{ESPRIT} pour
l'encadrement académique et les connaissances transmises durant mon cursus.

Un remerciement tout particulier à \textbf{M.~Harrathi Hafedh}, mon encadrant
académique, pour son accompagnement, ses orientations pertinentes et le temps
qu'il m'a consacré tout au long de ce projet de fin d'études.

Je suis particulièrement reconnaissant envers \textbf{ma famille} pour leur
soutien indéfectible, leur patience et leurs encouragements tout au long de
mes études.

Un grand merci à \textbf{mes amis} et camarades de promotion pour leur
entraide, les échanges fructueux et les moments de partage qui ont enrichi
cette expérience académique.

% ────────────────────────────────────────────────────────────
% RÉSUMÉ
% ────────────────────────────────────────────────────────────
\chapter*{Résumé}
\addcontentsline{toc}{chapter}{Résumé}

Ce rapport présente la conception et le développement de \textbf{LogiWay},
une plateforme web complète dédiée à la gestion intelligente de flottes
de véhicules poids lourds avec traçabilité GPS en temps réel.
Le projet a été réalisé chez \textbf{ExypnoTech Engineering Services}
à Monastir, Tunisie.

La plateforme adopte une architecture multi-couches : un backend
\textbf{Spring Boot 3.x} sécurisé par \textbf{Keycloak} (OAuth2/JWT),
un frontend \textbf{Angular 17}, une base de données \textbf{MySQL 8},
et trois microservices Python d'intelligence artificielle.

\medskip
Les modules réalisés couvrent :
\begin{itemize}
  \item La \textbf{sécurité et l'authentification} via Keycloak (OIDC, JWT, RBAC)
  \item La \textbf{gestion multi-entreprises} avec workflow de validation
  \item La \textbf{gestion des véhicules, secteurs et utilisateurs}
  \item Le \textbf{suivi GPS temps réel} via WebSocket STOMP et OSRM
  \item La \textbf{gestion des congés et réclamations} avec validation IA
  \item Un \textbf{système de notifications SSE} persistées en MySQL
  \item La \textbf{messagerie interne} en temps réel (WebSocket)
  \item Un \textbf{module ML de pauses réglementaires} (RandomForest, CE 561/2006)
  \item Un \textbf{assistant chatbot RAG} (LangChain, Google Gemini, FastAPI)
  \item Des \textbf{tableaux de bord KPI} analytiques personnalisés par rôle
  \item Une \textbf{stratégie de tests complète} (JUnit 5, pytest, Jest, k6, Playwright)
\end{itemize}

\medskip
\textbf{Mots-clés :} Spring Boot, Keycloak, JWT, OAuth2, Angular 17,
MySQL, JPA/Hibernate, WebSocket, SSE, OSRM, Leaflet, GNSS, RandomForest,
RAG, LangChain, Gemini, Scrum, Flask, FastAPI, JUnit 5, Playwright, k6.

\newpage

% ────────────────────────────────────────────────────────────
% ABSTRACT
% ────────────────────────────────────────────────────────────
\chapter*{Abstract}
\addcontentsline{toc}{chapter}{Abstract}

This report presents the design and development of \textbf{LogiWay}, a
comprehensive web platform for intelligent management of heavy-vehicle fleets
with real-time GPS tracking. The project was carried out at
\textbf{ExypnoTech Engineering Services}, Monastir, Tunisia.

The platform combines a \textbf{Spring Boot 3.x} backend secured with
\textbf{Keycloak} (OAuth2/JWT), an \textbf{Angular 17} frontend, a
\textbf{MySQL 8} database, and three Python AI microservices. The implemented
modules cover: authentication, multi-company management, fleet and vehicle
management, geographic sectors, GPS trip planning and tracking,
leave and complaint management with AI validation, real-time SSE notifications,
internal messaging, ML-based regulatory break detection (RandomForest, EU
Regulation 561/2006), a RAG chatbot (LangChain + Gemini), role-based KPI
dashboards, and a complete automated testing strategy.

\medskip
\textbf{Keywords:} Fleet management, Spring Boot, Keycloak, Angular,
WebSocket, GPS, OSRM, Machine Learning, RandomForest, RAG, LangChain,
Google Gemini, Scrum, AI microservices, JUnit 5, Playwright, k6.

\newpage

% ────────────────────────────────────────────────────────────
% ABRÉVIATIONS
% ────────────────────────────────────────────────────────────
\chapter*{Liste des abréviations}
\addcontentsline{toc}{chapter}{Liste des abréviations}

\begin{tabularx}{\textwidth}{lX}
  \textbf{API}     & Application Programming Interface \\
  \textbf{CRUD}    & Create, Read, Update, Delete \\
  \textbf{DTO}     & Data Transfer Object \\
  \textbf{E2E}     & End-to-End (tests de bout en bout) \\
  \textbf{GNSS}    & Global Navigation Satellite System \\
  \textbf{GPS}     & Global Positioning System \\
  \textbf{HTTP}    & HyperText Transfer Protocol \\
  \textbf{IAM}     & Identity and Access Management \\
  \textbf{JWT}     & JSON Web Token \\
  \textbf{KPI}     & Key Performance Indicator \\
  \textbf{LLM}     & Large Language Model \\
  \textbf{MAE}     & Mean Absolute Error \\
  \textbf{ML}      & Machine Learning \\
  \textbf{MVC}     & Model-View-Controller \\
  \textbf{NLU}     & Natural Language Understanding \\
  \textbf{OIDC}    & OpenID Connect \\
  \textbf{OSRM}    & Open Source Routing Machine \\
  \textbf{RAG}     & Retrieval-Augmented Generation \\
  \textbf{RBAC}    & Role-Based Access Control \\
  \textbf{REST}    & Representational State Transfer \\
  \textbf{RF}      & RandomForest \\
  \textbf{ROPC}    & Resource Owner Password Credentials \\
  \textbf{SPA}     & Single Page Application \\
  \textbf{SQL}     & Structured Query Language \\
  \textbf{SSE}     & Server-Sent Events \\
  \textbf{STOMP}   & Simple Text Oriented Messaging Protocol \\
  \textbf{SMTP}    & Simple Mail Transfer Protocol \\
  \textbf{UML}     & Unified Modeling Language \\
  \textbf{VU}      & Virtual User (utilisateur virtuel de test) \\
  \textbf{VRP}     & Vehicle Routing Problem \\
\end{tabularx}

\newpage

% ────────────────────────────────────────────────────────────
% TABLES
% ────────────────────────────────────────────────────────────
\tableofcontents\newpage
\listoffigures\newpage
\listoftables\newpage

\clearpage
\pagenumbering{arabic}
\setcounter{page}{1}
\pagestyle{fancy}


% ============================================================
% CHAPITRE 1 — PRÉSENTATION GÉNÉRALE
% ============================================================
\chapitreavecplan{Présentation générale du projet}
\begin{plandechapitre}
  \planitem{Introduction}{sec:ch1-introduction}
  \planitem{Présentation de l'organisme d'accueil}{sec:ch1-organisme}
  \planitem{Contexte du secteur du transport routier tunisien}{sec:ch1-contexte}
  \planitem{Problématique}{sec:ch1-problematique}
  \planitem{Étude de l'existant}{sec:ch1-existant}
  \planitem{Solution proposée : LogiWay}{sec:ch1-solution}
  \planitem{Méthodologie : Agile Scrum et organisation en 6 sprints}{sec:ch1-methodologie}
\end{plandechapitre}

\section{Introduction}
\label{sec:ch1-introduction}

Ce premier chapitre présente l'organisme d'accueil, le contexte du secteur
du transport routier tunisien, la problématique identifiée, les solutions
existantes sur le marché, la solution proposée LogiWay et la méthodologie
Agile Scrum retenue pour piloter le développement en 6 sprints.

\section{Présentation de l'organisme d'accueil}
\label{sec:ch1-organisme}

\subsection{ExypnoTech Engineering Services}

ExypnoTech Engineering Services, fondée en mars 2023 par
\textbf{M.~Wajih El~Hadj Youssef} (PDG) et \textbf{M.~Fehmi Kharroubi}
(Directeur Technique), est une entreprise innovante basée à Monastir, Tunisie.
Initialement spécialisée dans les solutions avancées pour l'aquaculture,
la société élargit aujourd'hui son champ d'action à divers projets
d'ingénierie logicielle.

\begin{figure}[H]
  \centering
  \includegraphics[width=0.35\textwidth]{log_exy.jpg}
  \caption{Logo ExypnoTech Engineering Services}
  \label{fig:logo_exy}
\end{figure}

En tant que participant actif du programme \textbf{ACTincube}, dirigé par
Actia Engineering Services, ExypnoTech bénéficie d'un mentorat de 24 mois.

\subsection{Domaines d'activité}

ExypnoTech intervient dans trois domaines complémentaires :
\begin{itemize}
  \item \textbf{Solutions numériques innovantes} : plateformes web et mobiles,
        systèmes embarqués, traitement de données en temps réel.
  \item \textbf{Engagement pour la durabilité} : pratiques responsables,
        sécurité et résilience opérationnelle.
  \item \textbf{Impact global} : solutions à visée internationale.
\end{itemize}

\section{Contexte du secteur du transport routier tunisien}
\label{sec:ch1-contexte}

Le transport routier constitue \textbf{l'épine dorsale de l'économie
tunisienne}, assurant plus de 80\% du fret domestique. Pourtant, la majorité
des opérateurs continuent de gérer leurs flottes au moyen d'outils rudimentaires.

Cette situation génère des pertes opérationnelles considérables :
retards non détectés, surcoûts carburant, accidents liés à la fatigue,
et incapacité à produire des analyses de performance.

Le cadre réglementaire européen CE~561/2006, appliqué aux transporteurs
tunisiens travaillant avec des donneurs d'ordre européens, impose des
contraintes strictes. Une infraction peut entraîner des amendes allant
jusqu'à \textbf{15~000~€}.

\section{Problématique}
\label{sec:ch1-problematique}

\begin{tcolorbox}[keybox]
  \textbf{Problématique centrale :} Comment concevoir et développer une
  plateforme numérique intégrée permettant aux opérateurs de flottes
  tunisiennes de gérer en temps réel leurs véhicules, chauffeurs et missions,
  tout en garantissant la conformité réglementaire, la sécurité des conducteurs
  et la traçabilité complète des opérations~-- à coût maîtrisé et sans
  dépendance à des solutions propriétaires étrangères~?
\end{tcolorbox}

\medskip
Six problèmes majeurs ont été identifiés :

\begin{enumerate}
  \item \textbf{Absence de traçabilité GPS adaptée} aux contraintes poids lourds.
  \item \textbf{Gestion dispersée et manuelle} sans centralisation.
  \item \textbf{Non-respect des pauses réglementaires} : risques d'amendes.
  \item \textbf{Absence d'analyse prédictive} pour anticiper les retards.
  \item \textbf{Communication défaillante} entre direction, managers et chauffeurs.
  \item \textbf{Empreinte carbone non maîtrisée}.
\end{enumerate}

\section{Étude de l'existant}
\label{sec:ch1-existant}

\subsection{Analyse des solutions du marché}

\begin{table}[H]
  \centering
  \caption{Benchmark des solutions de gestion de flotte existantes}
  \label{tab:benchmark}
  \rowcolors{2}{tablerow}{white}
  \begin{tabularx}{\textwidth}{|l|Y|Y|}
    \hline
    \thead{Solution} & \thead{Points forts} & \thead{Limites principales} \\
    \hline
    \textbf{Webfleet / TomTom} &
      GPS temps réel, trafic live &
      26~€/véhicule/mois, peu adapté au contexte tunisien \\
    \textbf{Geotab} &
      Télématique OBD avancée, API ouverte &
      Tarification sur devis, intégration complexe \\
    \textbf{Teksat} &
      Géolocalisation flottes françaises &
      Pas de module HR, pas de chatbot IA \\
    \textbf{Samsara} &
      Caméras IoT, alertes comportementales &
      Abonnement très élevé, données hébergées aux USA \\
    \textbf{FleetComplete} &
      Optimisation itinéraires &
      Interface complexe, API fermée \\
    \hline
  \end{tabularx}
\end{table}

\subsection{Analyse SWOT}

\begin{table}[H]
  \centering
  \caption{Analyse SWOT de la solution LogiWay}
  \label{tab:swot}
  \begin{tabularx}{\textwidth}{|Y|Y|}
    \hline
    \multicolumn{2}{|c|}{\cellcolor{tablehead}\textcolor{white}{\textbf{Facteurs internes}}} \\
    \hline
    \cellcolor{green!10}\textbf{Forces} &
    \cellcolor{red!10}\textbf{Faiblesses} \\
    \hline
    \begin{itemize}
      \item Solution open-source, auto-hébergée
      \item Sécurité OAuth2/JWT centralisée
      \item IA intégrée (ML, RAG, expert system)
      \item Architecture multi-rôles évolutive
    \end{itemize}
    &
    \begin{itemize}
      \item Nécessite une connexion réseau stable
      \item Dépendance Keycloak pour l'IAM
      \item Services Flask/FastAPI IA séparés
    \end{itemize}
    \\
    \hline
    \multicolumn{2}{|c|}{\cellcolor{tablehead}\textcolor{white}{\textbf{Facteurs externes}}} \\
    \hline
    \cellcolor{blue!10}\textbf{Opportunités} &
    \cellcolor{orange!10}\textbf{Menaces} \\
    \hline
    \begin{itemize}
      \item Fort besoin numérique du transport TN
      \item Réduction des coûts opérationnels
      \item Expansion des modules IA (Phase~2)
    \end{itemize}
    &
    \begin{itemize}
      \item Concurrence solutions établies
      \item Évolution des réglementations
      \item Dépendance API externe Gemini
    \end{itemize}
    \\
    \hline
  \end{tabularx}
\end{table}

\section{Solution proposée : LogiWay}
\label{sec:ch1-solution}

LogiWay répond à l'ensemble des lacunes identifiées en proposant :

\begin{itemize}
  \item \textbf{Plateforme open-source, auto-hébergée} : coûts maîtrisés.
  \item \textbf{Authentification Keycloak OAuth2/JWT} avec trois rôles
        (SuperAdmin, Manager, Chauffeur) et cookies HttpOnly.
  \item \textbf{Gestion multi-entreprises} avec isolation stricte des données.
  \item \textbf{Suivi GPS temps réel} Multi-GNSS avec OSRM poids lourds.
  \item \textbf{Gestion RH complète} : congés, réclamations, messagerie.
  \item \textbf{Trois modules IA} : validation réclamations (système expert),
        pauses réglementaires (RandomForest ML), chatbot RAG (Gemini + LangChain).
  \item \textbf{Notifications SSE et WebSocket} temps réel.
  \item \textbf{Tableaux de bord KPI} personnalisés par rôle.
  \item \textbf{Tests automatisés complets} : unitaires, performance, E2E.
\end{itemize}

\section{Méthodologie : Agile Scrum et organisation en 6 sprints}
\label{sec:ch1-methodologie}

\subsection{Justification du choix Agile Scrum}

La méthodologie \textbf{Agile Scrum} a été choisie pour :
\begin{itemize}
  \item \textbf{Itérations courtes} (2 semaines par sprint) permettant des
        validations fréquentes avec l'encadrant.
  \item \textbf{Backlog priorisé} : modules critiques développés en premier.
  \item \textbf{Livraisons progressives} : chaque sprint produit un incrément
        fonctionnel testable.
  \item \textbf{Adaptabilité} aux spécifications évolutives du projet.
\end{itemize}

\subsection{Organisation en 6 sprints}

Le projet LogiWay est organisé en \textbf{6 sprints de 2 semaines chacun},
soit 12 semaines de développement. Chaque sprint produit un incrément
fonctionnel livrable suivant la dépendance naturelle des fonctionnalités.

\begin{table}[H]
  \centering
  \caption{Organisation des 6 sprints du projet LogiWay}
  \label{tab:sprints}
  \rowcolors{2}{tablerow}{white}
  \begin{tabularx}{0.97\textwidth}{|c|Y|c|c|}
    \hline
    \thead{Sprint} & \thead{Thème principal} & \thead{Durée} & \thead{Priorité} \\
    \hline
    \textbf{Sprint 1} &
      Fondations : Authentification + Utilisateurs + Entreprises + Secteurs + Véhicules &
      2 semaines & CRITIQUE \\
    \textbf{Sprint 2} &
      Opérations RH : Gestion des Congés + Système de Notifications SSE &
      2 semaines & HAUTE \\
    \textbf{Sprint 3} &
      Communication : Messagerie interne + Chatbot IA RAG &
      2 semaines & MOYENNE \\
    \textbf{Sprint 4} &
      C\oe{}ur métier : Trajets + Pauses ML (CE 561/2006) + Réclamations IA &
      2 semaines & CRITIQUE \\
    \textbf{Sprint 5} &
      Intelligence : Tableaux de bord KPI + Analytique &
      2 semaines & HAUTE \\
    \textbf{Sprint 6} &
      Qualité \& Déploiement : Tests automatisés + DevOps Docker &
      2 semaines & CRITIQUE \\
    \hline
  \end{tabularx}
\end{table}

\subsection{Rôles dans la plateforme LogiWay}

\begin{table}[H]
  \centering
  \caption{Rôles et périmètres des acteurs de LogiWay}
  \label{tab:roles}
  \rowcolors{2}{tablerow}{white}
  \begin{tabularx}{\textwidth}{|l|Y|}
    \hline
    \thead{Rôle} & \thead{Périmètre et responsabilités} \\
    \hline
    \textbf{SUPERADMIN} &
      Vision globale multi-entreprises, validation des comptes, modération
      des réclamations urgentes, accès à toutes les analytics \\
    \textbf{MANAGER} &
      Gestion de son entreprise : flotte, chauffeurs, trajets, secteurs,
      réclamations de son équipe, validation des congés \\
    \textbf{DRIVER (Chauffeur)} &
      Espace personnel : ses trajets, alertes pauses ML en temps réel,
      ses congés, ses réclamations, cockpit de conduite \\
    \hline
  \end{tabularx}
\end{table}

\section{Conclusion}

Ce chapitre a posé le cadre général du projet LogiWay : le contexte du
transport routier tunisien, les lacunes des solutions existantes, la solution
proposée et la méthodologie Agile Scrum adoptée en 6 sprints.
Les chapitres suivants détaillent l'architecture technique puis chacun
des 6 sprints réalisés.

\newpage


% ============================================================
% CHAPITRE 2 — ÉTUDE TECHNIQUE ET ARCHITECTURE
% ============================================================
\chapitreavecplan{Étude technique et architecture}
\begin{plandechapitre}
  \planitem{Introduction}{sec:ch2-introduction}
  \planitem{Architecture globale de LogiWay}{sec:ch2-archi-globale}
  \planitem{Architecture logique}{sec:ch2-archi-logique}
  \planitem{Architecture physique}{sec:ch2-archi-physique}
  \planitem{Comparaison entre l'architecture logique et physique}{sec:ch2-comparaison}
  \planitem{Technologies utilisées}{sec:ch2-technologies}
  \planitem{Outils et services d'infrastructure}{sec:ch2-outils}
  \planitem{Modèle de données}{sec:ch2-modele-donnees}
  \planitem{Diagramme de classes}{sec:ch2-diagramme-classes}
  \planitem{Synthèse de l'architecture}{sec:ch2-synthese}
\end{plandechapitre}

\section{Introduction}
\label{sec:ch2-introduction}

Ce chapitre présente l'étude technique et l'architecture globale de la
plateforme LogiWay. LogiWay repose sur une architecture logicielle
principalement monolithique basée sur Spring Boot, complétée par Angular,
Keycloak, MySQL et plusieurs services Python d'intelligence artificielle.

\section{Architecture globale de LogiWay}
\label{sec:ch2-archi-globale}

La plateforme LogiWay adopte une architecture globale composée de :
\begin{itemize}
    \item \textbf{Frontend Angular} : interface utilisateur SPA ;
    \item \textbf{Backend Spring Boot} : c\oe{}ur métier et API REST ;
    \item \textbf{Keycloak} : authentification et autorisation ;
    \item \textbf{MySQL} : stockage persistant des données ;
    \item \textbf{Services IA et ML} : fonctionnalités intelligentes ;
    \item \textbf{OSRM} : calcul des itinéraires via API publique HTTPS ;
    \item \textbf{Leaflet et OpenStreetMap} : cartographie interactive ;
    \item \textbf{Simulateur GPS Python} : génération de positions GPS.
\end{itemize}

\section{Architecture logique}
\label{sec:ch2-archi-logique}

\subsection{Architecture logique du frontend}

L'application Angular est organisée selon :
\begin{itemize}
    \item \textbf{Components (UI)} : affichage des interfaces et interactions ;
    \item \textbf{Services HTTP} : communication avec les API REST ;
    \item \textbf{Guards} : protection des routes selon l'authentification ;
    \item \textbf{Interceptor JWT} : ajout automatique du token dans les requêtes ;
    \item \textbf{Models (TypeScript)} : représentation des données échangées.
\end{itemize}

\subsection{Architecture logique du backend}

Le backend Spring Boot adopte une organisation en couches :
\begin{itemize}
    \item \textbf{Couche de présentation} : contrôleurs REST exposant les API ;
    \item \textbf{Couche métier} : services implémentant les règles fonctionnelles ;
    \item \textbf{Couche d'accès aux données} : repositories Spring Data JPA ;
    \item \textbf{DTO et Mapper} : contrôle des données échangées entre couches ;
    \item \textbf{Sécurité} : Spring Security intégré avec Keycloak.
\end{itemize}

\begin{figure}[H]
    \centering
    \includegraphics[width=1\textwidth]{architecture_logique_final_corrige.png}
    \caption{Architecture logique de la plateforme LogiWay}
    \label{fig:architecture_logique}
\end{figure}

\section{Architecture physique}
\label{sec:ch2-archi-physique}

\begin{table}[H]
  \centering
  \caption{Déploiement physique des composants de LogiWay}
  \label{tab:archi_physique}
  \rowcolors{2}{tablerow}{white}
  \begin{tabularx}{\textwidth}{|l|c|Y|}
    \hline
    \thead{Composant} & \thead{Port} & \thead{Description} \\
    \hline
    Frontend Angular & 4200 & Interface utilisateur Web \\
    Backend Spring Boot & 8080 & API REST principale et logique métier \\
    Keycloak & 8180 & Authentification et gestion des identités \\
    MySQL & 3306 & Base de données relationnelle \\
    Service ML Pauses & 5000 & Flask + RandomForest (pauses CE 561/2006) \\
    Service IA Réclamations & 5001 & Flask + système expert (validation) \\
    Service RAG Chatbot & 5003 & FastAPI + LangChain + Gemini \\
    OSRM Public API & 443 HTTPS & Calcul d'itinéraires via API publique \\
    \hline
  \end{tabularx}
\end{table}

\begin{figure}[H]
    \centering
    \includegraphics[width=1\textwidth]{architecture_pysique_final_corrige.png}
    \caption{Architecture physique de la plateforme LogiWay}
    \label{fig:architecture_physique}
\end{figure}

\section{Comparaison entre l'architecture logique et physique}
\label{sec:ch2-comparaison}

\begin{table}[H]
  \centering
  \caption{Comparaison entre les deux architectures}
  \label{tab:comparaison_architectures}
  \rowcolors{2}{tablerow}{white}
  \begin{tabularx}{\textwidth}{|l|Y|Y|}
    \hline
    \thead{Élément} &
    \thead{Architecture logique} &
    \thead{Architecture physique} \\
    \hline
    Objectif &
    Organisation interne du logiciel &
    Déploiement réel du système \\
    Niveau &
    Fonctionnel et logiciel &
    Infrastructure et réseau \\
    Principaux éléments &
    Couches, composants, services &
    Serveurs, ports, processus \\
    Exemple &
    Controller $\rightarrow$ Service $\rightarrow$ Repository &
    Angular:4200 $\rightarrow$ Spring Boot:8080 \\
    \hline
  \end{tabularx}
\end{table}

\section{Technologies utilisées}
\label{sec:ch2-technologies}

\begin{table}[H]
    \centering
    \caption{Synthèse des technologies et outils utilisés dans LogiWay}
    \label{tab:technologies_logiway}
    \rowcolors{2}{tablerow}{white}
    \begin{tabularx}{\textwidth}{|l|l|Y|}
        \hline
        \thead{Catégorie} & \thead{Technologie / Outil} & \thead{Utilisation} \\
        \hline
        Frontend & Angular 17+ & Application web SPA \\
        Frontend & Kendo UI & Composants d'interface \\
        Frontend & Leaflet & Cartographie interactive \\
        Frontend & TypeScript & Langage principal côté client \\
        Backend & Java 21 / Spring Boot 3.x & API REST et logique métier \\
        Backend & Spring Data JPA / Hibernate & Persistance des données \\
        Backend & Spring Security & Sécurisation des API \\
        IA & Python / Flask & Services ML et IA \\
        IA & FastAPI & Service chatbot RAG \\
        IA/ML & scikit-learn & Modèle Random Forest \\
        IA & Google Gemini & Génération de réponses IA \\
        Routage & OSRM & Calcul des itinéraires \\
        Base de données & MySQL 8 & Stockage des données métier \\
        Authentification & Keycloak 24 & Gestion des identités et rôles \\
        Versioning & Git / GitHub & Gestion du code source \\
        Build & Maven / npm & Gestion des dépendances \\
        Déploiement & Docker & Conteneurisation \\
        \hline
    \end{tabularx}
\end{table}

\begin{figure}[H]
    \centering
    \includegraphics[width=1\textwidth]{outils_technologie.png}
    \caption{Vue d'ensemble des technologies et outils utilisés dans LogiWay}
    \label{fig:technologies_logiway}
\end{figure}

\section{Outils et services d'infrastructure}
\label{sec:ch2-outils}

\begin{table}[H]
  \centering
  \caption{Outils et services utilisés dans le projet}
  \label{tab:outils_services}
  \rowcolors{2}{tablerow}{white}
  \begin{tabularx}{\textwidth}{|l|Y|}
    \hline
    \thead{Outil / Service} & \thead{Utilisation} \\
    \hline
    Git / GitHub & Gestion du code source et collaboration \\
    Maven & Gestion des dépendances backend Spring Boot \\
    Docker & Conteneurisation et simplification du déploiement \\
    Swagger / OpenAPI & Documentation et test des API REST \\
    Leaflet / OpenStreetMap & Affichage cartographique \\
    OSRM Public API & Calcul des itinéraires via API publique HTTPS \\
    \hline
  \end{tabularx}
\end{table}

\section{Modèle de données}
\label{sec:ch2-modele-donnees}

\begin{table}[H]
  \centering
  \caption{Entités principales du modèle de données LogiWay}
  \label{tab:entites}
  \rowcolors{2}{tablerow}{white}
  \begin{tabularx}{\textwidth}{|l|Y|}
    \hline
    \thead{Entité} & \thead{Relations principales} \\
    \hline
    \textbf{Utilisateur} & Classe de base pour tous les utilisateurs \\
    \textbf{SuperAdmin} & Gestion globale de la plateforme \\
    \textbf{Manager} & Gestion des chauffeurs, véhicules, missions \\
    \textbf{Entreprise} & Gestion des informations de l'entreprise \\
    \textbf{Chauffeur} & Association avec Manager et Véhicule \\
    \textbf{Vehicule} & Association avec Chauffeur et Trajet \\
    \textbf{Trajet} & Gestion des déplacements et itinéraires \\
    \textbf{PauseReglementaire} & Association avec les trajets \\
    \textbf{Conge} & Gestion des congés des chauffeurs \\
    \textbf{Reclamation} & Gestion et traitement des réclamations \\
    \textbf{Notification} & Notifications destinées aux utilisateurs \\
    \textbf{MessengerMessage} & Échanges de messages entre utilisateurs \\
    \hline
  \end{tabularx}
\end{table}

\section{Diagramme de classes}
\label{sec:ch2-diagramme-classes}

\begin{figure}[H]
    \centering
    \includegraphics[width=1\textwidth]{classe1.png}
    \caption{Diagramme de classes UML -- principales entités de LogiWay}
    \label{fig:class_diagram}
\end{figure}

\section{Synthèse de l'architecture}
\label{sec:ch2-synthese}

L'architecture de LogiWay repose sur une organisation centralisée autour
du backend monolithique Spring Boot. Le frontend Angular constitue le point
d'interaction principal avec les utilisateurs. Keycloak assure la sécurité.
MySQL assure la persistance. Les services IA Python sont déployés séparément
et communiquent via HTTP/REST.

\section{Conclusion}

Ce chapitre a présenté l'étude technique et l'architecture de LogiWay.
L'architecture logique décrit l'organisation interne des composants,
l'architecture physique le déploiement réel. Les technologies Angular,
Spring Boot, MySQL, Keycloak, Python, Flask, FastAPI et OSRM constituent
la base technique sur laquelle reposent les 6 sprints de développement.

\newpage


% ============================================================
% CHAPITRE 3 — SPRINT 1 : FONDATIONS
% ============================================================
\chapitreavecplan{Sprint 1 --- Fondations : Authentification, Utilisateurs, Entreprises, Secteurs et Véhicules}
\begin{plandechapitre}
  \planitem{Introduction}{sec:ch3-introduction}
  \planitem{Authentification et gestion des utilisateurs}{sec:ch3-sprint1}
  \planitem{Gestion des entreprises}{sec:ch3-sprint2}
  \planitem{Gestion des véhicules et de la flotte}{sec:ch3-sprint3}
  \planitem{Gestion des secteurs géographiques}{sec:ch3-sprint4}
  \planitem{Conclusion du Sprint 1}{sec:ch3-conclusion}
\end{plandechapitre}

\section{Introduction}
\label{sec:ch3-introduction}

Le Sprint~1 met en place la \textbf{fondation complète} de la plateforme
LogiWay : authentification sécurisée, gestion des rôles, structure
organisationnelle (entreprises, secteurs, véhicules). Ce sprint est le
préalable indispensable à tous les autres, car rien ne peut fonctionner
sans les utilisateurs, les entreprises et les véhicules.

\textbf{Rôles concernés :}
\begin{itemize}
    \item SUPERADMIN : gère toutes les entreprises, valide les comptes, crée les managers
    \item MANAGER : gère son équipe, sa flotte, ses secteurs
    \item DRIVER : se connecte, consulte son profil et son véhicule assigné
\end{itemize}

\section{Authentification et gestion des utilisateurs}
\label{sec:ch3-sprint1}

\subsection{Présentation du module}

Le premier module est consacré au système d'authentification et de gestion
des utilisateurs. La solution repose sur Keycloak pour la gestion de
l'identité et des rôles, tandis que Spring Boot assure la sécurisation
des ressources et la gestion des règles métier.

\subsection{Flux d'authentification Keycloak}

Le flux d'authentification repose sur une architecture OAuth2/OIDC intégrant
Keycloak comme fournisseur d'identité. L'utilisateur saisit ses identifiants
depuis Angular. Une fois l'authentification réussie, un token JWT est délivré.
Spring Boot joue le rôle de \textbf{Resource Server OAuth2} et vérifie la
validité du JWT à partir des clés publiques exposées par Keycloak.

Les rôles sont récupérés depuis le token pour appliquer le contrôle d'accès
(RBAC). Après login, la redirection est automatique selon le rôle :
SUPERADMIN~$\rightarrow$~\texttt{/dashboard/superadmin},
MANAGER~$\rightarrow$~\texttt{/dashboard/manager},
DRIVER~$\rightarrow$~\texttt{/dashboard/driver}.

\begin{figure}[H]
  \centering
  \includegraphics[width=0.98\textwidth]{diagramme_sequence_authentification.png}
  \caption{Diagramme de séquence du processus d'authentification avec Keycloak}
  \label{fig:seq_auth}
\end{figure}

\begin{figure}[H]
  \centering
  \includegraphics[width=0.95\textwidth]{diagramme-activite-athentification.png}
  \caption{Diagramme d'activité du processus d'authentification}
  \label{fig:activite_auth}
\end{figure}

\subsection{Fonctionnalités d'authentification}

\begin{table}[H]
  \centering
  \caption{Fonctionnalités d'authentification et de gestion des comptes}
  \label{tab:auth_features}
  \rowcolors{2}{tablerow}{white}
  \begin{tabularx}{\textwidth}{|l|l|Y|}
    \hline
    \thead{Fonctionnalité} & \thead{Rôle} & \thead{Description} \\
    \hline
    Login / Logout & Tous & Authentification via Keycloak OAuth2/OIDC \\
    Forgot password & Tous & Code de réinitialisation par email \\
    Reset password & Tous & Validation du code et mise à jour \\
    Vérification email & Manager & Confirmation lors de la création du compte \\
    Refresh token & Tous & Renouvellement automatique du token \\
    Créer Manager & SuperAdmin & Création et synchronisation MySQL/Keycloak \\
    Créer Chauffeur & SuperAdmin/Manager & Création et assignation au Manager \\
    Modifier utilisateur & SuperAdmin & Mise à jour des informations \\
    Activer/Désactiver & SuperAdmin & Gestion du statut d'accès \\
    \hline
  \end{tabularx}
\end{table}

\begin{figure}[H]
  \centering
  \includegraphics[width=1\textwidth]{3d518737-fcef-4988-a684-fe9790ca3aeb.png}
  \caption{Flux complet du processus d'authentification : Angular, Keycloak et Spring Boot}
  \label{fig:auth_flow}
\end{figure}

\section{Gestion des entreprises}
\label{sec:ch3-sprint2}

\subsection{Présentation du module}

Ce module permet au Manager de créer et gérer son entreprise. Le SuperAdmin
assure la validation avant d'autoriser l'accès complet aux modules
opérationnels.

\subsection{Workflow de validation}

\[
\texttt{EN\_ATTENTE}
\rightarrow
\texttt{ACTIF}
\quad\text{ou}\quad
\texttt{EN\_ATTENTE}
\rightarrow
\texttt{REJETE}
\]

\begin{figure}[H]
  \centering
  \includegraphics[width=1\textwidth]{sequence_Créer une Entreprise.png}
  \caption{Diagramme de séquence de création et validation d'une entreprise}
  \label{fig:seq_entreprise}
\end{figure}

\begin{figure}[H]
  \centering
  \includegraphics[width=0.75\textwidth]{diagramme_activite_vie entreprise.png}
  \caption{Diagramme d'activité du workflow d'une entreprise}
  \label{fig:activite_entreprise}
\end{figure}

\subsection{Règles métier}

\begin{itemize}
  \item Un Manager ne peut posséder qu'une seule entreprise.
  \item La validation, le rejet ou la suspension déclenche une notification
        et un email aux utilisateurs concernés.
  \item Un Manager dont l'entreprise est \texttt{SUSPENDU} n'accède plus
        aux modules opérationnels.
\end{itemize}

\section{Gestion des véhicules et de la flotte}
\label{sec:ch3-sprint3}

\subsection{Présentation du module}

Ce module permet aux Managers de gérer les véhicules de leur entreprise
et d'assurer leur affectation aux chauffeurs. Le backend vérifie
systématiquement que le chauffeur et le véhicule appartiennent à la même
entreprise (\texttt{HTTP 409 Conflict} en cas d'incohérence).

\begin{figure}[H]
  \centering
  \includegraphics[width=1\textwidth]{diagramme_Sequence_Planification_Assignation_chauffeur_vehicul.png}
  \caption{Diagramme de séquence d'affectation d'un véhicule à un chauffeur}
  \label{fig:seq_vehicule}
\end{figure}

\begin{figure}[H]
  \centering
  \includegraphics[width=1\textwidth]{activite_vehicule1.png}
  \caption{Diagramme d'activité de gestion d'un véhicule}
  \label{fig:activite_vehicule}
\end{figure}

\begin{table}[H]
  \centering
  \caption{Entité Véhicule -- attributs principaux}
  \label{tab:vehicule}
  \rowcolors{2}{tablerow}{white}
  \begin{tabularx}{\textwidth}{|l|l|Y|}
    \hline
    \thead{Attribut} & \thead{Type} & \thead{Description} \\
    \hline
    matricule & String (UNIQUE) & Immatriculation unique \\
    marque / modele & String & Marque et modèle \\
    capacite & Double & Capacité de charge (tonnes) \\
    kilometrage & Integer & Kilométrage total \\
    statut & Enum & EN\_SERVICE, EN\_MAINTENANCE, HORS\_SERVICE \\
    chauffeurActuel & Chauffeur (FK) & Chauffeur actuellement affecté \\
    entreprise & Entreprise (FK) & Entreprise propriétaire \\
    \hline
  \end{tabularx}
\end{table}

\section{Gestion des secteurs géographiques}
\label{sec:ch3-sprint4}

\subsection{Présentation du module}

Un secteur représente une zone géographique opérationnelle pour organiser
les activités de l'entreprise et faciliter l'affectation des Managers.

\begin{figure}[H]
  \centering
  \includegraphics[width=1\textwidth]{sequene_Créer un Secteur.png}
  \caption{Diagramme de séquence d'assignation d'un Manager à un secteur}
  \label{fig:seq_secteur}
\end{figure}

\begin{figure}[H]
  \centering
  \includegraphics[width=1\textwidth]{activite_secteur.png}
  \caption{Diagramme d'activité de gestion des secteurs géographiques}
  \label{fig:activite_secteur}
\end{figure}

\subsection{Règles métier}

\begin{itemize}
  \item Un secteur est associé à une seule entreprise.
  \item Une zone géographique ne peut pas être dupliquée.
  \item Un Manager ne peut être assigné qu'à un seul secteur.
  \item L'assignation déclenche une notification au Manager concerné.
\end{itemize}

\section{Conclusion du Sprint 1}
\label{sec:ch3-conclusion}

Le Sprint~1 établit la fondation complète de LogiWay. L'authentification
Keycloak, la gestion multi-entreprises avec workflow de validation,
la gestion de la flotte avec isolation par entreprise, et la structuration
géographique en secteurs constituent le socle sur lequel reposent tous
les sprints suivants.

\newpage


% ============================================================
% CHAPITRE 4 — SPRINT 2 : CONGÉS + NOTIFICATIONS
% ============================================================
\chapitreavecplan{Sprint 2 --- Opérations RH : Gestion des Congés et Système de Notifications}
\begin{plandechapitre}
  \planitem{Introduction}{sec:ch4-introduction}
  \planitem{Gestion des congés}{sec:ch4-conges}
  \planitem{Système de notifications SSE}{sec:ch4-sse}
  \planitem{Conclusion du Sprint 2}{sec:ch4-conclusion}
\end{plandechapitre}

\section{Introduction}
\label{sec:ch4-introduction}

Le Sprint~2 développe les opérations RH quotidiennes de la plateforme :
demandes de congé des chauffeurs et système de notifications en temps réel.
Les notifications SSE constituent une infrastructure partagée utilisée par
tous les sprints suivants pour alerter les acteurs des événements métier.

\textbf{Rôles concernés :}
\begin{itemize}
    \item DRIVER : soumet des demandes de congé, reçoit les notifications personnelles
    \item MANAGER : approuve ou refuse les congés, reçoit les alertes de son équipe
    \item SUPERADMIN : vue globale des congés, reçoit les notifications critiques
\end{itemize}

\section{Gestion des congés}
\label{sec:ch4-conges}

\subsection{Présentation du module}

Le module de gestion des congés permet aux chauffeurs de soumettre des
demandes de congé et aux managers de les valider ou refuser. Un calendrier
visuel permet de visualiser les absences planifiées de l'équipe.

\subsection{Workflow de validation}

\[
\texttt{EN\_ATTENTE}
\rightarrow
\texttt{APPROUVE}
\quad\text{ou}\quad
\texttt{EN\_ATTENTE}
\rightarrow
\texttt{REJETE}
\]

Les types de congé disponibles : \texttt{MALADIE}, \texttt{MARIAGE},
\texttt{VACANCES}. Un chauffeur en statut \texttt{EN\_SERVICE}
ne peut pas soumettre une nouvelle demande.

\begin{figure}[H]
  \centering
  \includegraphics[width=0.90\textwidth]{diagramme_Sequence_Demande_Validation_Conge2.png}
  \caption{Diagramme de séquence de demande et validation de congé}
  \label{fig:seq_conge}
\end{figure}

\begin{figure}[H]
  \centering
  \includegraphics[width=0.98\textwidth]{diagramme_Activite_Circuit_Gestion_Conges.png}
  \caption{Diagramme d'activité du workflow de validation d'un congé}
  \label{fig:activite_conge}
\end{figure}

\section{Système de notifications SSE}
\label{sec:ch4-sse}

\subsection{Présentation du module}

Le système de notifications repose sur la technologie
\textbf{Server-Sent Events (SSE)} qui permet au serveur d'envoyer des
événements aux clients via une connexion HTTP persistante.

\begin{tcolorbox}[alertbox]
  \textbf{Règle absolue :} Toute notification est d'abord persistée
  dans MySQL avant d'être diffusée via SSE. Si l'utilisateur n'est pas
  connecté, la notification reste disponible en base de données.
\end{tcolorbox}

\begin{figure}[H]
  \centering
  \includegraphics[width=\textwidth]{diagramme_Sequence_gestion_Notification.png}
  \caption{Diagramme de séquence de diffusion d'une notification SSE}
  \label{fig:seq_sse}
\end{figure}

\begin{figure}[H]
  \centering
  \includegraphics[width=1\textwidth]{diagramme_Activite_Traitement_Alerte.png}
  \caption{Diagramme d'activité du traitement d'une notification SSE}
  \label{fig:activite_sse}
\end{figure}

\subsection{Événements déclencheurs par rôle}

\begin{table}[H]
  \centering
  \caption{Événements SSE, destinataires et types}
  \label{tab:sse_events}
  \rowcolors{2}{tablerow}{white}
  \begin{tabularx}{\textwidth}{|Y|l|l|}
    \hline
    \thead{Événement} & \thead{Destinataire} & \thead{Type} \\
    \hline
    Nouvelle entreprise EN\_ATTENTE & SuperAdmins & NOTIF\_ENTREPRISE \\
    Validation / rejet entreprise & Manager propriétaire & NOTIF\_ENTREPRISE \\
    Trajet planifié et assigné & Chauffeur concerné & NOTIF\_TRAJET \\
    Demande de congé soumise & Manager du chauffeur & NOTIF\_CONGE \\
    Congé approuvé ou rejeté & Chauffeur demandeur & NOTIF\_CONGE \\
    Réclamation urgente & SuperAdmins + Manager & NOTIF\_RECLAMATION \\
    Affectation véhicule & Chauffeur concerné & NOTIF\_VEHICULE \\
    Assignation Manager à secteur & Manager assigné & NOTIF\_SECTEUR \\
    Pauses IA générées & Manager du trajet & NOTIF\_TRAJET \\
    \hline
  \end{tabularx}
\end{table}

\section{Conclusion du Sprint 2}
\label{sec:ch4-conclusion}

Le Sprint~2 fournit la couche de communication événementielle de LogiWay.
Le système SSE est une infrastructure transversale réutilisée par les
Sprints~3 à~5 pour diffuser les alertes de pauses, les notifications de
réclamations et les messages internes.

\newpage

% ============================================================
% CHAPITRE 5 — SPRINT 3 : MESSAGERIE + CHATBOT RAG
% ============================================================
\chapitreavecplan{Sprint 3 --- Communication : Messagerie Interne et Chatbot Intelligent RAG}
\begin{plandechapitre}
  \planitem{Introduction}{sec:ch5-introduction}
  \planitem{Messagerie interne WebSocket}{sec:ch5-messenger}
  \planitem{Chatbot RAG et génération de rapports}{sec:ch5-rag}
  \planitem{Technologies du Sprint 3}{sec:ch5-technologies}
  \planitem{Conclusion du Sprint 3}{sec:ch5-conclusion}
\end{plandechapitre}

\section{Introduction}
\label{sec:ch5-introduction}

Le Sprint~3 apporte les outils de communication de la plateforme :
une messagerie directe en temps réel entre utilisateurs et un assistant
IA conversationnel basé sur l'approche RAG (Retrieval-Augmented Generation).

\textbf{Rôles concernés :}
\begin{itemize}
    \item SUPERADMIN et MANAGER : accès complet au chatbot RAG et à la génération de rapports
    \item DRIVER : messagerie avec son manager, consultation du chatbot sur ses trajets personnels
\end{itemize}

\section{Messagerie interne WebSocket}
\label{sec:ch5-messenger}

\subsection{Présentation du module}

La messagerie interne repose sur \textbf{WebSocket STOMP} pour la
communication temps réel. Les messages et l'historique sont persistés
dans MySQL. Les règles d'accès par rôle contrôlent qui peut contacter qui :
SUPERADMIN avec tous, MANAGER avec ses chauffeurs, DRIVER avec son manager.

\subsection{Fonctionnalités}

\begin{itemize}
  \item Conversations privées 1-to-1 persistées dans MySQL.
  \item Historique complet des échanges.
  \item Gestion des statuts de lecture (vu/non vu).
  \item Compteur de messages non lus dans la barre de navigation.
  \item Notification SSE à réception d'un nouveau message (Sprint 2).
\end{itemize}

\begin{figure}[H]
  \centering
  \includegraphics[width=1\textwidth]{diagramme_Sequence_Envoyer_Message_Messenger.png}
  \caption{Diagramme de séquence d'envoi et réception d'un message}
  \label{fig:seq_messagerie}
\end{figure}

\begin{figure}[H]
  \centering
  \includegraphics[width=1\textwidth]{diagramme_ Activité — traitement_messenger.png}
  \caption{Diagramme d'activité du fonctionnement de la messagerie interne}
  \label{fig:activite_messagerie}
\end{figure}

\section{Chatbot intelligent RAG et génération de rapports}
\label{sec:ch5-rag}

\subsection{Architecture globale}

\begin{table}[H]
  \centering
  \caption{Couches de l'architecture chatbot RAG}
  \label{tab:archi_rag}
  \rowcolors{2}{tablerow}{white}
  \begin{tabularx}{\textwidth}{|l|Y|}
    \hline
    \thead{Couche} & \thead{Rôle} \\
    \hline
    Frontend Angular &
      Interface conversationnelle, formulaire de génération, téléchargement automatique \\
    Backend Spring Boot &
      Authentification JWT, contrôle RBAC, proxy sécurisé vers FastAPI \\
    Service RAG FastAPI (port 5003) &
      LangChain + Google Gemini : NLU, génération SQL, création PDF/CSV/TXT \\
    MySQL &
      Données métier, historique des conversations \\
    \hline
  \end{tabularx}
\end{table}

\subsection{Module 1 : Chatbot intelligent avec RAG}

Le RAG (\textit{Retrieval-Augmented Generation}) combine :
\begin{itemize}
  \item \textbf{Retrieval} : extraction des données pertinentes depuis MySQL en temps réel.
  \item \textbf{Generation} : Google Gemini formule une réponse en français naturel.
\end{itemize}

Flux de traitement :
\begin{enumerate}
  \item L'utilisateur pose une question en français via Angular.
  \item Spring Boot valide le JWT et les droits RBAC.
  \item Gemini analyse l'intention (domaine, action, entités).
  \item LangChain sélectionne l'outil approprié.
  \item Requête SQL exécutée sur MySQL.
  \item Gemini génère la réponse avec le contexte des données réelles.
\end{enumerate}

\begin{table}[H]
  \centering
  \caption{Outils LangChain du chatbot LogiWay}
  \label{tab:langchain_tools}
  \rowcolors{2}{tablerow}{white}
  \begin{tabularx}{\textwidth}{|l|Y|}
    \hline
    \thead{Outil} & \thead{Fonction} \\
    \hline
    Obtenir informations véhicules & Statut, disponibilité et détails de la flotte \\
    Obtenir informations chauffeurs & Données chauffeurs, affectations, disponibilité \\
    Obtenir informations trajets & Historique et statut des trajets \\
    Obtenir informations congés & Gestion et statut des congés \\
    Obtenir informations réclamations & Suivi des réclamations \\
    Exécuter requête SQL & Requêtes dynamiques sécurisées \\
    \hline
  \end{tabularx}
\end{table}

\subsection{Module 2 : Génération intelligente de rapports}

La génération transforme une requête en langage naturel en document
professionnel :

\[
  \text{Requête NL} \xrightarrow{\text{NLU}} \text{Paramètres}
  \xrightarrow{\text{SQL}} \text{Données}
  \xrightarrow{\text{Générateur}} \text{PDF/CSV/TXT}
\]

\begin{table}[H]
  \centering
  \caption{Paramètres extraits par l'analyse NLU}
  \label{tab:nlu}
  \rowcolors{2}{tablerow}{white}
  \begin{tabularx}{\textwidth}{|l|l|Y|}
    \hline
    \thead{Paramètre} & \thead{Valeurs} & \thead{Exemple} \\
    \hline
    Domaine & véhicules, chauffeurs, trajets, congés, réclamations &
      ``liste des véhicules'' $\rightarrow$ domaine = véhicules \\
    Format & PDF, CSV, TXT &
      ``rapport en CSV'' $\rightarrow$ format = CSV \\
    Période & aujourd'hui, semaine, mois, année &
      ``ce mois'' $\rightarrow$ date\_debut = premier du mois \\
    Filtres & statut, priorité, tri, limite &
      ``réclamations ouvertes'' $\rightarrow$ statut = EN\_COURS \\
    \hline
  \end{tabularx}
\end{table}

\section{Technologies du Sprint 3}
\label{sec:ch5-technologies}

\begin{table}[H]
  \centering
  \caption{Technologies utilisées -- messagerie et chatbot RAG}
  \label{tab:techno_chatbot}
  \rowcolors{2}{tablerow}{white}
  \begin{tabularx}{\textwidth}{|l|Y|}
    \hline
    \thead{Technologie} & \thead{Rôle} \\
    \hline
    WebSocket STOMP & Messagerie temps réel \\
    Google Gemini & LLM pour compréhension NL et génération \\
    LangChain & Orchestration des agents et outils métier \\
    FastAPI & Service web Python hautes performances \\
    SQLAlchemy & ORM Python, requêtes paramétrées sécurisées \\
    ReportLab & Génération PDF professionnels \\
    \hline
  \end{tabularx}
\end{table}

\section{Conclusion du Sprint 3}
\label{sec:ch5-conclusion}

Le Sprint~3 enrichit LogiWay de deux outils de communication complémentaires.
La messagerie WebSocket couvre les besoins immédiats de coordination.
Le chatbot RAG transforme l'accès à l'information en permettant des
requêtes en français naturel sur toutes les données de la plateforme.

\newpage

% ============================================================
% CHAPITRE 6 — SPRINT 4 : CŒUR MÉTIER
% ============================================================
\chapitreavecplan{Sprint 4 --- C\oe{}ur Métier : Trajets GPS, Pauses Réglementaires ML et Réclamations IA}
\begin{plandechapitre}
  \planitem{Introduction}{sec:ch6-introduction}
  \planitem{Gestion des trajets et géolocalisation GPS}{sec:ch6-trajets}
  \planitem{Module ML : Pauses réglementaires CE 561/2006}{sec:ch6-pauses}
  \planitem{Module IA : Validation intelligente des réclamations}{sec:ch6-reclamations}
  \planitem{Conclusion du Sprint 4}{sec:ch6-conclusion}
\end{plandechapitre}

\section{Introduction}
\label{sec:ch6-introduction}

Le Sprint~4 constitue le c\oe{}ur métier logistique de LogiWay : planification
et suivi GPS des trajets poids lourds, détection intelligente des pauses
réglementaires (CE~561/2006) par Machine Learning, et validation IA des
réclamations. Ce sprint est le plus complexe techniquement et nécessite
la base stable établie par les Sprints~1 et~2.

\textbf{Rôles concernés :}
\begin{itemize}
    \item DRIVER : conduit avec alertes pauses ML en direct, soumet des réclamations validées par IA
    \item MANAGER : planifie les trajets, supervise les pauses, traite les réclamations de son équipe
    \item SUPERADMIN : analytics globales de conformité, modère les réclamations urgentes
\end{itemize}

\section{Gestion des trajets et géolocalisation GPS}
\label{sec:ch6-trajets}

\subsection{Systèmes GNSS intégrés}

\begin{table}[H]
  \centering
  \caption{Systèmes GNSS intégrés et avantages pour LogiWay}
  \label{tab:gnss}
  \rowcolors{2}{tablerow}{white}
  \begin{tabularx}{\textwidth}{|l|l|c|Y|}
    \hline
    \thead{Système} & \thead{Opérateur} & \thead{Précision} & \thead{Avantages} \\
    \hline
    GPS     & USA              & 3--5 m  & Couverture mondiale, fiabilité maximale \\
    Galileo & Union Européenne & 1 m     & Meilleure précision civile \\
    GLONASS & Russie           & 3--5 m  & Performant en relief montagneux \\
    BeiDou  & Chine            & 5--10 m & Bonne couverture Afrique du Nord \\
    \hline
  \end{tabularx}
\end{table}

\subsection{Calcul d'itinéraire OSRM}

OSRM calcule des itinéraires adaptés aux poids lourds :
\begin{enumerate}
  \item Le Manager saisit départ, destination et arrêts intermédiaires.
  \item Spring Boot appelle l'API OSRM avec le profil \texttt{truck}.
  \item OSRM retourne la géométrie GeoJSON, la distance et la durée.
  \item Le backend stocke la géométrie en MySQL.
  \item Angular affiche le tracé Leaflet avec marqueurs colorés.
\end{enumerate}

\subsection{Suivi GPS temps réel}

\begin{itemize}
  \item Le Chauffeur démarre le trajet, Angular active la géolocalisation.
  \item Chaque position GPS est envoyée au backend via WebSocket.
  \item Le backend la persiste et la diffuse au Manager en temps réel.
  \item Si l'écart dépasse \textbf{500~m}, une alerte est envoyée et
        un recalcul OSRM est déclenché automatiquement.
\end{itemize}

\begin{table}[H]
  \centering
  \caption{Tables du module Trajets et GPS}
  \label{tab:trajets_db}
  \rowcolors{2}{tablerow}{white}
  \begin{tabularx}{\textwidth}{|l|Y|}
    \hline
    \thead{Table} & \thead{Colonnes principales} \\
    \hline
    \texttt{trajets} &
      id, pointDepart, destination, dateDepart, statut (EN\_COURS/ACTIF/COMPLETE),
      chauffeurId, vehiculeId, managerId, geometrieItineraire (JSON) \\
    \texttt{positions\_gps} &
      id, trajetId, latitude, longitude, vitesse, cap, timestamp \\
    \texttt{pauses\_reglementaires} &
      id, trajetId, type, latitude, longitude, aiScore, fatigueScore, reasoning \\
    \hline
  \end{tabularx}
\end{table}

\section{Module ML : Pauses réglementaires CE 561/2006}
\label{sec:ch6-pauses}

\subsection{Contexte et problématique}

La fatigue au volant représente \textbf{30\% des accidents mortels} impliquant
des poids lourds. Le Règlement Européen CE~561/2006 impose un encadrement
strict des temps de conduite avec des sanctions allant jusqu'à \textbf{15~000~€
par infraction}.

\subsection{Seuils réglementaires}

\begin{table}[H]
  \centering
  \caption{Seuils réglementaires appliqués par le système}
  \label{tab:reglementation}
  \rowcolors{2}{tablerow}{white}
  \begin{tabularx}{\textwidth}{|l|Y|Y|}
    \hline
    \thead{Seuil} & \thead{Règle} & \thead{Action système} \\
    \hline
    \textbf{3 heures} &
      Alerte préventive recommandée &
      WARNING\_ALERT : popup non bloquante chauffeur + notification Manager \\
    \textbf{4h30} &
      Pause de 45 min OBLIGATOIRE &
      MANDATORY\_REST : popup non fermable + alerte critique SuperAdmin \\
    \textbf{9h/jour} &
      Maximum journalier de conduite &
      Enregistrement conformité, rapport audit \\
    \hline
  \end{tabularx}
\end{table}

\subsection{Calcul du temps de conduite}

\begin{tcolorbox}[ruleframe]
\[
  T_{\text{conduite}} = \frac{T_{\text{écoulé}} - \sum T_{\text{pauses}}}{3600}
  \text{ heures}
\]
Un scheduler Spring Boot évalue tous les trajets EN\_COURS toutes les
\textbf{2 minutes}, filtrant ceux avec $T_{\text{conduite}} \geq 3$ heures.
\end{tcolorbox}

\subsection{Le modèle RandomForest}

\begin{table}[H]
  \centering
  \caption{Justification du choix de RandomForest}
  \label{tab:rf_justif}
  \rowcolors{2}{tablerow}{white}
  \begin{tabularx}{\textwidth}{|l|Y|}
    \hline
    \thead{Critère} & \thead{Avantage RandomForest} \\
    \hline
    Précision & R² = 0.90 (90\% de précision prédictive) \\
    Robustesse & Résiste au bruit et aux valeurs aberrantes \\
    Performance & Prédiction en $< 50$ ms \\
    Généralisation & Pas de sur-apprentissage grâce à l'ensemble d'arbres \\
    \hline
  \end{tabularx}
\end{table}

\begin{figure}[H]
  \centering
  \includegraphics[width=1\textwidth]{pause_ai1.png}
  \caption{Flux complet du processus Machine Learning pour les pauses}
  \label{fig:auth_pause_ml}
\end{figure}

\subsection{Les 15 features analysées}

\begin{table}[H]
  \centering
  \caption{Features du modèle RandomForest -- groupes et poids}
  \label{tab:features}
  \rowcolors{2}{tablerow}{white}
  \begin{tabularx}{\textwidth}{|l|l|Y|c|}
    \hline
    \thead{Groupe} & \thead{Feature} & \thead{Description} & \thead{Poids} \\
    \hline
    \multirow{3}{*}{État chauffeur} &
      Heures de conduite & Heures continues & 35\% \\
    & Distance totale & Distance du trajet en km & 12\% \\
    & Progression & Ratio avancement trajet (0--1) & 18\% \\
    \hline
    \multirow{2}{*}{Contexte temporel} &
      Heure d'arrivée & Heure estimée au POI & 10\% \\
    & Heure de repas & Indicateur heure repas & 8\% \\
    \hline
    \multirow{5}{*}{Caractéristiques POI} &
      Distance POI--route & Distance perpendiculaire & 15\% \\
    & Accès poids lourds & Autorisation HGV & 25\% \\
    & Douches & Présence de douches & 5\% \\
    & Toilettes & Présence de toilettes & 8\% \\
    & Ouvert 24h & Disponibilité permanente & 7\% \\
    \hline
  \end{tabularx}
\end{table}

\subsection{Interprétation des scores IA}

\begin{table}[H]
  \centering
  \caption{Interprétation des scores de recommandation}
  \label{tab:scores_ia}
  \rowcolors{2}{tablerow}{white}
  \begin{tabularx}{\textwidth}{|c|l|Y|}
    \hline
    \thead{Score} & \thead{Niveau} & \thead{Interprétation} \\
    \hline
    0--40   & Non recommandé & POI inadapté (trop loin, pas d'accès PL) \\
    41--69  & Acceptable     & POI convenable mais non optimal \\
    70--84  & Recommandé     & Bonne combinaison timing + accessibilité PL \\
    85--100 & Urgent         & Seuil 4h30 proche -- conformité CE 561/2006 \\
    \hline
  \end{tabularx}
\end{table}

\section{Module IA : Validation intelligente des réclamations}
\label{sec:ch6-reclamations}

\subsection{Architecture du système}

\begin{table}[H]
  \centering
  \caption{Couches du système de validation des réclamations}
  \label{tab:archi_recla}
  \rowcolors{2}{tablerow}{white}
  \begin{tabularx}{\textwidth}{|l|Y|}
    \hline
    \thead{Couche} & \thead{Rôle} \\
    \hline
    Frontend Angular &
      Formulaire réactif, debounce 600~ms, validation temps réel \\
    Backend Spring Boot &
      Appel du service IA avant persistance via RestTemplate \\
    Service IA Python Flask (port 5001) &
      Système expert : toxicité + sémantique. JSON : \{valide, typeErreur, scores\} \\
    \hline
  \end{tabularx}
\end{table}

\begin{figure}[H]
  \centering
  \includegraphics[width=1.1\textwidth]{flux_ai_reclamation.png}
  \caption{Flux de validation IA d'une réclamation -- de Angular à MySQL}
  \label{fig:recla_ia}
\end{figure}

\begin{tcolorbox}[keybox]
  \textbf{Important :} Le service de validation est un
  \textbf{système expert à règles déterministes} (Rule-Based Expert System),
  non un modèle Machine Learning. Il utilise des expressions régulières
  avec word boundaries pour la détection.
\end{tcolorbox}

\subsection{Couche 1 : Détection de toxicité}

Le dictionnaire contient plus de 30 mots en français et en anglais.
Seuil de blocage : $S_{\text{toxicité}} \geq 0.55$.

\begin{itemize}
  \item ``Ce \textbf{connard} ne répond pas'' $\rightarrow$ DÉTECTÉ (mot entier)
  \item ``Re\textbf{connaître} le problème'' $\rightarrow$ IGNORÉ (sous-chaîne)
\end{itemize}

\subsection{Couche 2 : Validation sémantique métier}

Seuil de validation : $S_{\text{sémantique}} \geq 0.25$ (minimum 2 mots métier).

\begin{table}[H]
  \centering
  \caption{Dictionnaire métier de validation sémantique}
  \label{tab:dictionnaire}
  \rowcolors{2}{tablerow}{white}
  \begin{tabularx}{\textwidth}{|l|Y|c|}
    \hline
    \thead{Domaine} & \thead{Exemples de termes} & \thead{Nb} \\
    \hline
    Véhicules & véhicule, camion, maintenance, panne, frein, moteur & 15 \\
    Trajets & trajet, route, itinéraire, livraison, secteur, retard & 18 \\
    Personnel & chauffeur, conducteur, pause, repos, horaire & 8 \\
    Incidents & incident, accident, problème, entrepôt, client & 9 \\
    \hline
  \end{tabularx}
\end{table}

\subsection{Logique de décision}

\[
  \text{Décision} =
  \begin{cases}
    \text{REJET (toxique)} & \text{si } S_{\text{toxicité}} \geq 0.55 \\
    \text{REJET (hors-sujet)} & \text{si } S_{\text{sémantique}} < 0.25 \\
    \text{ACCEPTATION} & \text{sinon}
  \end{cases}
\]

\begin{table}[H]
  \centering
  \caption{Exemples de résultats de validation}
  \label{tab:exemples_valid}
  \rowcolors{2}{tablerow}{white}
  \begin{tabularx}{\textwidth}{|Y|c|c|c|}
    \hline
    \thead{Texte soumis} & \thead{$S_{tox}$} & \thead{$S_{sem}$} & \thead{Résultat} \\
    \hline
    ``Le camion 42 a une panne de frein urgente'' &
      0.0 & 1.0 & \textcolor{teal}{\textbf{ACCEPTÉ}} \\
    ``Ce connard ne répond jamais'' &
      1.0 & --- & \textcolor{red}{\textbf{REJETÉ (toxique)}} \\
    ``J'ai mal à la tête aujourd'hui'' &
      0.0 & 0.0 & \textcolor{amber}{\textbf{REJETÉ (hors-sujet)}} \\
    \hline
  \end{tabularx}
\end{table}

\begin{figure}[H]
  \centering
  \includegraphics[width=0.90\textwidth]{diagramme_Séquence —Soumettre une Réclamation1.png}
  \caption{Diagramme de séquence de soumission d'une réclamation}
  \label{fig:seq_reclamation}
\end{figure}

\begin{figure}[H]
  \centering
  \includegraphics[width=1\textwidth]{diagramme_Activité 1— Traitement d'une Réclamation.png}
  \caption{Diagramme d'activité du traitement d'une réclamation}
  \label{fig:activite_reclamation}
\end{figure}

\section{Conclusion du Sprint 4}
\label{sec:ch6-conclusion}

Le Sprint~4 livre les fonctionnalités différenciantes de LogiWay. La
combinaison Multi-GNSS + OSRM poids lourds + Leaflet assure un suivi
précis. Le modèle RandomForest automatise la conformité CE~561/2006.
Le système expert de validation des réclamations garantit la qualité
du contenu soumis.

\newpage


% ============================================================
% CHAPITRE 7 — SPRINT 5 : DASHBOARD KPIs ET ANALYTIQUE
% ============================================================
\chapitreavecplan{Sprint 5 --- Intelligence : Tableaux de Bord KPI et Analytique}
\begin{plandechapitre}
  \planitem{Introduction}{sec:ch7-introduction}
  \planitem{Dashboard SUPERADMIN -- Vue globale}{sec:ch7-superadmin}
  \planitem{Dashboard MANAGER -- Console de gestion}{sec:ch7-manager}
  \planitem{Dashboard DRIVER -- Espace Cockpit}{sec:ch7-driver}
  \planitem{KPIs du module Pauses Réglementaires}{sec:ch7-pauses-kpi}
  \planitem{Graphiques analytiques (7 visualisations)}{sec:ch7-graphiques}
  \planitem{Classement des chauffeurs -- Top 10}{sec:ch7-classement}
  \planitem{Filtres et statistiques avancées}{sec:ch7-filtres}
  \planitem{Conclusion du Sprint 5}{sec:ch7-conclusion}
\end{plandechapitre}

\section{Introduction}
\label{sec:ch7-introduction}

Le Sprint~5 offre une vision globale et analytique de toute la plateforme
via des tableaux de bord KPI adaptés à chaque rôle. Ce sprint est développé
en dernier parmi les sprints fonctionnels car il agrège les données de tous
les modules précédents : utilisateurs, véhicules, trajets, pauses, congés
et réclamations.

LogiWay propose des tableaux de bord analytiques personnalisés par rôle,
offrant une vision 360° des opérations. L'approche s'inspire des principes
Power~BI : KPIs clés, graphiques dynamiques, filtres interactifs.

\textbf{Accès et périmètre :}

\begin{table}[H]
  \centering
  \caption{Accès aux tableaux de bord par rôle}
  \label{tab:dashboard_roles}
  \rowcolors{2}{tablerow}{white}
  \begin{tabularx}{\textwidth}{|l|Y|}
    \hline
    \thead{Rôle} & \thead{Périmètre des données} \\
    \hline
    SUPERADMIN & Toutes les entreprises, tous les chauffeurs, toutes les flottes \\
    MANAGER & Uniquement les données de son entreprise \\
    DRIVER & Données personnelles uniquement (cockpit de conduite) \\
    \hline
  \end{tabularx}
\end{table}

\section{Dashboard SUPERADMIN -- Vue globale}
\label{sec:ch7-superadmin}

\subsection{KPI Cards -- 8 indicateurs principaux}

Les KPI Cards sont affichées en haut du dashboard. Elles permettent une
lecture rapide de l'état opérationnel global.

\begin{table}[H]
  \centering
  \caption{KPI Cards du dashboard SUPERADMIN}
  \label{tab:kpi_superadmin}
  \rowcolors{2}{tablerow}{white}
  \begin{tabularx}{\textwidth}{|l|Y|Y|}
    \hline
    \thead{KPI} & \thead{Description} & \thead{Utilité} \\
    \hline
    Chauffeurs actifs &
      Nombre EN\_SERVICE / total inscrit &
      Détecter les indisponibilités imprévues \\
    Véhicules en service &
      EN\_SERVICE / EN\_MAINTENANCE / HORS\_SERVICE &
      Suivre les immobilisations de la flotte \\
    Missions en cours &
      Trajets actifs (badge retard si applicable) &
      Identifier les livraisons problématiques \\
    Taux de ponctualité &
      (missions à l'heure / total) × 100 &
      Mesurer la qualité de service globale \\
    Congés en attente &
      Demandes non traitées (alerte si $> 5$) &
      Éviter le retard de traitement RH \\
    Réclamations ouvertes &
      EN\_COURS / RÉSOLU / total &
      Mesurer l'efficacité de modération \\
    Utilisateurs actifs &
      SUPERADMIN / MANAGER / DRIVER &
      Suivi de la composition de la plateforme \\
    Global Score Flotte &
      Score pondéré sur 100 points &
      Synthèse unique de la santé opérationnelle \\
    \hline
  \end{tabularx}
\end{table}

\subsection{Score Global Flotte -- Calcul pondéré}

\begin{tcolorbox}[ruleframe]
\[
  S_{\text{global}} = 0.35 \cdot S_{\text{ponctualité}}
  + 0.25 \cdot S_{\text{véhicules}}
  + 0.20 \cdot S_{\text{missions}}
  + 0.10 \cdot S_{\text{carburant}}
  + 0.10 \cdot S_{\text{incidents}}
\]
\end{tcolorbox}

\begin{figure}[H]
  \centering
  \begin{subfigure}[b]{0.48\textwidth}
    \centering
    \includegraphics[width=\textwidth]{dhaboard_superadmin.png}
    \caption{Vue d'ensemble KPI Cards}
    \label{fig:superadmin_1}
  \end{subfigure}
  \hfill
  \begin{subfigure}[b]{0.48\textwidth}
    \centering
    \includegraphics[width=\textwidth]{dhaboard_superadmin1.png}
    \caption{Graphiques analytiques}
    \label{fig:superadmin_2}
  \end{subfigure}
  \caption{Dashboard SUPERADMIN -- Vue globale et analytique}
  \label{fig:superadmin}
\end{figure}

\section{Dashboard MANAGER -- Console de gestion}
\label{sec:ch7-manager}

KPIs filtrés exclusivement sur l'entreprise du Manager :
\begin{itemize}
  \item \textbf{Missions actives} : trajets EN\_COURS de l'entreprise
  \item \textbf{Véhicules disponibles} : en statut EN\_SERVICE
  \item \textbf{Historique 7 jours} : missions de la semaine
  \item \textbf{Secteur} : statut Activé / Non assigné
  \item \textbf{Taux de conformité pauses} : chauffeurs respectant CE~561/2006
  \item \textbf{Réclamations équipe} : statuts en cours / résolues
\end{itemize}

\section{Dashboard DRIVER -- Espace Cockpit}
\label{sec:ch7-driver}

KPIs personnels pour le chauffeur en cours de mission :
\begin{itemize}
  \item \textbf{Trajet en cours} : carte, itinéraire, ETA, distance restante
  \item \textbf{Compteur de conduite} : jauge visuelle avec décompte
  \item \textbf{Prochaine pause ML} : heure et durée prédites par RandomForest
  \item \textbf{Score de performance} : note sur 100 calculée en temps réel
  \item \textbf{Mes réclamations} : statut des réclamations en cours
  \item \textbf{Mes congés} : jours approuvés et à venir
\end{itemize}

\textbf{Graphique radar} sur 6 axes : Ponctualité, Sécurité, Consommation,
Pauses, Missions, Conformité.

\begin{figure}[H]
  \centering
  \begin{subfigure}[b]{0.48\textwidth}
    \centering
    \includegraphics[width=\textwidth]{dashboard_chauuffeur.png}
    \caption{Vue d'ensemble}
    \label{fig:chauffeur_1}
  \end{subfigure}
  \hfill
  \begin{subfigure}[b]{0.48\textwidth}
    \centering
    \includegraphics[width=\textwidth]{dashboard_chauuffeur1.png}
    \caption{Informations mission}
    \label{fig:chauffeur_2}
  \end{subfigure}
  \caption{Espace Chauffeur -- Cockpit de mission}
  \label{fig:chauffeur}
\end{figure}

\section{KPIs du module Pauses Réglementaires}
\label{sec:ch7-pauses-kpi}

\begin{table}[H]
  \centering
  \caption{KPIs du dashboard Pauses Réglementaires (Pause-AI)}
  \label{tab:kpi_pause}
  \rowcolors{2}{tablerow}{white}
  \begin{tabularx}{\textwidth}{|l|Y|}
    \hline
    \thead{KPI} & \thead{Description} \\
    \hline
    Pauses recommandées &
      Nombre d'alertes IA selon temps de conduite continu \\
    Pauses effectuées &
      Arrêts réellement validés par les chauffeurs \\
    Pauses ignorées &
      Recommandations ignorées (risque conformité CE 561/2006) \\
    Taux de conformité &
      $\frac{\text{Pauses effectuées}}{\text{Pauses recommandées}} \times 100$ \\
    \hline
  \end{tabularx}
\end{table}

\section{Graphiques analytiques -- 7 visualisations}
\label{sec:ch7-graphiques}

\begin{table}[H]
  \centering
  \caption{Graphiques du dashboard SUPERADMIN/MANAGER}
  \label{tab:graphiques}
  \rowcolors{2}{tablerow}{white}
  \begin{tabularx}{\textwidth}{|c|l|l|Y|}
    \hline
    \thead{\#} & \thead{Graphique} & \thead{Type Chart.js} & \thead{Objectif} \\
    \hline
    1 & Ponctualité 12 mois &
      Stacked Bar + Line &
      Analyser l'évolution mensuelle de la ponctualité sur 1 an \\
    2 & Flux des missions &
      Area Chart &
      Visualiser le volume d'activité et détecter les pics \\
    3 & Incidents du mois &
      Donut Chart &
      Répartir les incidents par catégorie (pannes, accidents, embouteillages) \\
    4 & Consommation carburant &
      Gauge (Doughnut 180°) &
      Mesurer l'efficacité énergétique vs cible \\
    5 & Capacité de charge &
      Donut Chart &
      Analyser l'utilisation de la capacité par type de véhicule \\
    6 & Répartition des rôles &
      Donut Chart &
      Visualiser la composition des utilisateurs \\
    7 & Congés par statut &
      Horizontal Bar &
      Mesurer le backlog RH et la politique d'approbation \\
    \hline
  \end{tabularx}
\end{table}

\section{Classement des chauffeurs -- Top 10}
\label{sec:ch7-classement}

\subsection{Score de performance chauffeur}

\begin{tcolorbox}[ruleframe]
\[
  S_{\text{perf}} = 0.35 \cdot S_{\text{ponct}}
  + 0.20 \cdot S_{\text{missions}}
  + 0.15 \cdot S_{\text{conduite}}
  + 0.15 \cdot S_{\text{dispo}}
  + 0.15 \cdot S_{\text{exp}}
  - 3 \cdot N_{\text{incidents}}
\]
\end{tcolorbox}

\begin{table}[H]
  \centering
  \caption{Badges de performance des chauffeurs}
  \label{tab:badges}
  \rowcolors{2}{tablerow}{white}
  \begin{tabularx}{\textwidth}{|c|l|Y|}
    \hline
    \thead{Score} & \thead{Badge} & \thead{Signification} \\
    \hline
    85--100 & \textbf{Excellent} & Chauffeur exemplaire, à valoriser \\
    70--84  & \textbf{Bon} & Performance satisfaisante \\
    50--69  & \textbf{Normal} & Dans la moyenne, suivi recommandé \\
    < 50    & \textbf{À améliorer} & Accompagnement nécessaire \\
    \hline
  \end{tabularx}
\end{table}

\section{Filtres et statistiques avancées}
\label{sec:ch7-filtres}

Tous les dashboards disposent de filtres dynamiques :
période (jour/semaine/mois/trimestre/année), entreprise (SUPERADMIN),
type de véhicule, statut chauffeur, zone géographique.

\textbf{Dashboards complémentaires :}
\begin{itemize}
  \item \textbf{Congés} : colonnes empilées validés/rejetés/en attente.
  \item \textbf{Flotte} : jauges santé mécanique, kilométrage par véhicule.
  \item \textbf{Réclamations} : donut par statut, délai moyen de traitement.
  \item \textbf{Historique trajets} : fréquence hebdomadaire, km quotidiens.
\end{itemize}

\section{Conclusion du Sprint 5}
\label{sec:ch7-conclusion}

Le Sprint~5 transforme les données opérationnelles brutes de LogiWay
en indicateurs actionnables. Les 8 KPI Cards, 7 graphiques analytiques
et le classement des chauffeurs offrent à chaque rôle une vision
complète et différenciée pour prendre des décisions éclairées.

\newpage

% ============================================================
% CHAPITRE 8 — SPRINT 6 : TESTS ET DEVOPS
% ============================================================
\chapitreavecplan{Sprint 6 --- Qualité et Déploiement : Tests Automatisés et DevOps}
\begin{plandechapitre}
  \planitem{Introduction}{sec:ch8-introduction}
  \planitem{Stratégie de tests -- Pyramide et objectifs}{sec:ch8-strategie}
  \planitem{Tests Unitaires Backend -- Spring Boot}{sec:ch8-unitaires-backend}
  \planitem{Tests Unitaires IA/ML -- Python}{sec:ch8-unitaires-python}
  \planitem{Tests Unitaires Frontend -- Angular}{sec:ch8-unitaires-frontend}
  \planitem{Tests de Performance -- k6}{sec:ch8-performance}
  \planitem{Tests End-to-End -- Playwright}{sec:ch8-e2e}
  \planitem{Containerisation et déploiement Docker}{sec:ch8-docker}
  \planitem{Bilan global des tests}{sec:ch8-bilan}
\end{plandechapitre}

\section{Introduction}
\label{sec:ch8-introduction}

Le Sprint~6 valide la qualité de l'ensemble de la plateforme par des tests
automatisés à tous les niveaux et prépare le déploiement en production via
Docker. Ce sprint est développé en dernier car les tests E2E et de
performance nécessitent une application complète et fonctionnelle.

\textbf{Environnement de test :}
\begin{itemize}
  \item Système d'exploitation : Windows 11, Intel Core i7, 16 Go RAM
  \item Stack testée : Spring Boot 3.2.2 · Angular 19 · Flask · FastAPI
        · MySQL 8.0 · Keycloak 24
  \item Date d'exécution : 07 septembre 2026
\end{itemize}

\section{Stratégie de tests -- Pyramide et objectifs}
\label{sec:ch8-strategie}

\subsection{La pyramide des tests LogiWay}

\begin{figure}[H]
\centering
\begin{tikzpicture}
  \fill[rougeEsprit!15] (0,0) -- (6,0) -- (3,4) -- cycle;
  \draw[rougeEsprit,thick] (0,0) -- (6,0) -- (3,4) -- cycle;
  \draw[rougeEsprit,dashed] (0.75,1) -- (5.25,1);
  \draw[rougeEsprit,dashed] (1.5,2) -- (4.5,2);
  \node at (3,0.5) {\small Tests Unitaires (JUnit/pytest/Jest) -- 45+ tests};
  \node at (3,1.5) {\small Performance (k6) -- 4 scénarios};
  \node at (3,2.7) {\small E2E (Playwright) -- 15 tests};
\end{tikzpicture}
\caption{Pyramide des tests du projet LogiWay}
\label{fig:pyramide_tests}
\end{figure}

\subsection{Résumé exécutif des tests}

\begin{table}[H]
  \centering
  \caption{Résumé des résultats de tests}
  \label{tab:resume_tests}
  \rowcolors{2}{tablerow}{white}
  \begin{tabularx}{\textwidth}{|l|l|c|c|c|}
    \hline
    \thead{Niveau} & \thead{Framework} & \thead{Tests} & \thead{Couverture} & \thead{Statut} \\
    \hline
    Unitaire Backend & JUnit 5 + Mockito + JaCoCo & 10 & 97,7\% lignes & \textcolor{teal}{✓ PASSÉ} \\
    Unitaire IA/ML & pytest + coverage.py & 15 & 99\% & \textcolor{teal}{✓ PASSÉ} \\
    Unitaire Frontend & Jest + Istanbul & 20+ & 100\% & \textcolor{teal}{✓ PASSÉ} \\
    Performance & k6 v2.2.0 & 4 scénarios & p(95)=38ms & \textcolor{teal}{✓ PASSÉ} \\
    End-to-End & Playwright 1.63.0 & 15 & --- & \textcolor{teal}{✓ PASSÉ} \\
    \hline
    \multicolumn{2}{|l|}{\textbf{Total}} & \textbf{60+} & & \textcolor{teal}{\textbf{✓ TOUS PASSÉS}} \\
    \hline
  \end{tabularx}
\end{table}

\section{Tests Unitaires Backend -- Spring Boot}
\label{sec:ch8-unitaires-backend}

\subsection{Objectif et frameworks}

Vérifier la logique métier du backend Java de manière isolée,
sans base de données ni réseau réel.

\begin{table}[H]
  \centering
  \caption{Frameworks des tests unitaires Backend}
  \label{tab:fw_backend}
  \rowcolors{2}{tablerow}{white}
  \begin{tabularx}{\textwidth}{|l|Y|}
    \hline
    \thead{Outil} & \thead{Rôle} \\
    \hline
    JUnit 5 & Framework principal d'exécution des tests Java \\
    Mockito & Simulation des dépendances (replace la BDD par des objets simulés) \\
    AssertJ & Bibliothèque d'assertions fluentes et lisibles \\
    JaCoCo 0.8.11 & Génération des rapports de couverture de code \\
    Spring Boot Starter Test 3.2.2 & Conteneur intégrant tous les outils \\
    \hline
  \end{tabularx}
\end{table}

\subsection{Couverture de code (JaCoCo)}

\begin{table}[H]
  \centering
  \caption{Résultats de couverture JaCoCo}
  \label{tab:jacoco}
  \rowcolors{2}{tablerow}{white}
  \begin{tabularx}{\textwidth}{|l|c|c|}
    \hline
    \thead{Métrique} & \thead{Résultat} & \thead{Détail} \\
    \hline
    Classes couvertes  & \textbf{100\%} & 171 / 171 \\
    Méthodes couvertes & \textbf{88,4\%} & 1~125 / 1~273 \\
    Branches couvertes & \textbf{91,8\%} & 3~834 / 4~175 \\
    Lignes couvertes   & \textbf{97,7\%} & 7~163 / 7~330 \\
    \hline
  \end{tabularx}
\end{table}

\subsection{Module Réclamations -- 4 scénarios testés}

\begin{table}[H]
  \centering
  \caption{Scénarios de tests du service Réclamation}
  \label{tab:tests_reclamation}
  \rowcolors{2}{tablerow}{white}
  \begin{tabularx}{\textwidth}{|c|Y|Y|c|}
    \hline
    \thead{\#} & \thead{Action testée} & \thead{Scénario} & \thead{Résultat} \\
    \hline
    1 & Consulter une réclamation & ID valide $\rightarrow$ données retournées & \textcolor{teal}{✓} \\
    2 & Consulter une réclamation inexistante & ID inconnu $\rightarrow$ exception levée & \textcolor{teal}{✓} \\
    3 & Résoudre une réclamation & EN\_COURS $\rightarrow$ RÉSOLU & \textcolor{teal}{✓} \\
    4 & Créer réclamation urgente & Priorité URGENT $\rightarrow$ notification SuperAdmin & \textcolor{teal}{✓} \\
    \hline
  \end{tabularx}
\end{table}

\subsection{Module Pauses Réglementaires -- 6 scénarios testés}

\begin{table}[H]
  \centering
  \caption{Scénarios de tests du service Pauses CE 561/2006}
  \label{tab:tests_pauses}
  \rowcolors{2}{tablerow}{white}
  \begin{tabularx}{\textwidth}{|c|Y|Y|c|}
    \hline
    \thead{\#} & \thead{Action testée} & \thead{Règle CE 561/2006} & \thead{Résultat} \\
    \hline
    1 & Récupérer pauses trajet court & Trajet < 3h $\rightarrow$ aucune pause & \textcolor{teal}{✓} \\
    2 & Vérifier alerte préventive & Trajet > 3h $\rightarrow$ WARNING\_ALERT présent & \textcolor{teal}{✓} \\
    3 & Vérifier pause obligatoire & Trajet > 4h30 $\rightarrow$ MANDATORY\_REST 45min & \textcolor{teal}{✓} \\
    4 & Test durée 3h exactement & 180 min $\rightarrow$ éligible aux pauses & \textcolor{teal}{✓} \\
    5 & Test durée 4h & 240 min $\rightarrow$ éligible & \textcolor{teal}{✓} \\
    6 & Test paramétré (3 durées) & 270, 360, 480 min $\rightarrow$ tous éligibles & \textcolor{teal}{✓} \\
    \hline
  \end{tabularx}
\end{table}

\section{Tests Unitaires IA/ML -- Python}
\label{sec:ch8-unitaires-python}

\subsection{Objectif et frameworks}

Valider la précision des algorithmes de traitement des réclamations.

\begin{table}[H]
  \centering
  \caption{Frameworks des tests unitaires Python IA/ML}
  \label{tab:fw_python}
  \rowcolors{2}{tablerow}{white}
  \begin{tabularx}{\textwidth}{|l|Y|}
    \hline
    \thead{Outil} & \thead{Rôle} \\
    \hline
    pytest & Framework de test Python, exécution et rapport \\
    pytest-cov & Calcul et affichage de la couverture de code \\
    pytest-flask & Client de test pour les endpoints Flask \\
    coverage.py 7.11.0 & Analyse détaillée de la couverture \\
    \hline
  \end{tabularx}
\end{table}

\subsection{Couverture (coverage.py)}

\begin{table}[H]
  \centering
  \caption{Résultats de couverture Python}
  \label{tab:coverage_python}
  \rowcolors{2}{tablerow}{white}
  \begin{tabularx}{\textwidth}{|l|c|c|}
    \hline
    \thead{Module} & \thead{Instructions} & \thead{Couverture} \\
    \hline
    Service Réclamation IA & 370 & \textbf{99\%} \\
    Modèle ML & 57 & \textbf{100\%} \\
    \textbf{Total} & \textbf{427} & \textbf{99\%} \\
    \hline
  \end{tabularx}
\end{table}

\subsection{Module Toxicité -- 7 scénarios}

\begin{table}[H]
  \centering
  \caption{Scénarios de tests de détection de toxicité}
  \label{tab:tests_toxicite}
  \rowcolors{2}{tablerow}{white}
  \begin{tabularx}{\textwidth}{|c|Y|Y|c|}
    \hline
    \thead{\#} & \thead{Action testée} & \thead{Comportement attendu} & \thead{Résultat} \\
    \hline
    1 & Analyser texte professionnel & Score toxicité = 0.0 & \textcolor{teal}{✓} \\
    2 & Analyser texte avec insulte & Score $\geq$ 0.55 (toxique détecté) & \textcolor{teal}{✓} \\
    3 & Analyser sous-chaîne ambiguë & Score = 0.0 (word boundary respecté) & \textcolor{teal}{✓} \\
    4 & Analyser texte en majuscules & Score $\geq$ 0.55 (insensible casse) & \textcolor{teal}{✓} \\
    5--7 & 3 textes logistiques (paramétré) & Score = 0.0 (aucun faux positif) & \textcolor{teal}{✓} \\
    \hline
  \end{tabularx}
\end{table}

\subsection{Module Sémantique -- 5 scénarios}

\begin{table}[H]
  \centering
  \caption{Scénarios de tests de validation sémantique}
  \label{tab:tests_semantique}
  \rowcolors{2}{tablerow}{white}
  \begin{tabularx}{\textwidth}{|c|Y|Y|c|}
    \hline
    \thead{\#} & \thead{Action testée} & \thead{Comportement attendu} & \thead{Résultat} \\
    \hline
    1 & Analyser texte hors-sujet & Score < 0.25 (rejeté) & \textcolor{teal}{✓} \\
    2 & Analyser texte véhicule/panne & Score $\geq$ 0.25 (accepté) & \textcolor{teal}{✓} \\
    3 & Analyser texte trajet/livraison & Score $\geq$ 0.25 (accepté) & \textcolor{teal}{✓} \\
    4 & Vérifier borne maximale & Score $\leq$ 1.0 (normalisé) & \textcolor{teal}{✓} \\
    5 & Vérifier seuil mots-clés & 2 mots > 1 mot $\rightarrow$ score supérieur & \textcolor{teal}{✓} \\
    \hline
  \end{tabularx}
\end{table}

\section{Tests Unitaires Frontend -- Angular}
\label{sec:ch8-unitaires-frontend}

\subsection{Objectif et frameworks}

Vérifier le comportement des composants Angular et des guards de navigation
de manière isolée, sans appels réseau réels.

\begin{table}[H]
  \centering
  \caption{Frameworks des tests unitaires Frontend}
  \label{tab:fw_frontend}
  \rowcolors{2}{tablerow}{white}
  \begin{tabularx}{\textwidth}{|l|Y|}
    \hline
    \thead{Outil} & \thead{Rôle} \\
    \hline
    Jest 29.7.0 & Framework de test JavaScript/TypeScript \\
    jest-preset-angular 16.0.0 & Adaptateur Jest pour Angular \\
    jest-environment-jsdom & Simulation du DOM navigateur \\
    Istanbul (intégré) & Rapport de couverture de code \\
    \hline
  \end{tabularx}
\end{table}

\subsection{Couverture (Istanbul/Jest)}

\begin{table}[H]
  \centering
  \caption{Résultats de couverture Frontend}
  \label{tab:coverage_frontend}
  \rowcolors{2}{tablerow}{white}
  \begin{tabularx}{\textwidth}{|l|c|c|}
    \hline
    \thead{Métrique} & \thead{Résultat} & \thead{Détail} \\
    \hline
    Statements & \textbf{100\%} & 20 / 20 \\
    Branches & \textbf{100\%} & 4 / 4 \\
    Functions & \textbf{100\%} & 4 / 4 \\
    Lines & \textbf{100\%} & 19 / 19 \\
    \hline
  \end{tabularx}
\end{table}

\subsection{Composant Authentification -- 4 scénarios}

\begin{table}[H]
  \centering
  \caption{Scénarios de tests du composant de connexion}
  \label{tab:tests_auth_front}
  \rowcolors{2}{tablerow}{white}
  \begin{tabularx}{\textwidth}{|c|Y|Y|c|}
    \hline
    \thead{\#} & \thead{Action testée} & \thead{Comportement attendu} & \thead{Résultat} \\
    \hline
    1 & Login valide & Redirection vers dashboard MANAGER & \textcolor{teal}{✓} \\
    2 & Mauvais mot de passe (erreur 401) & Message ``Invalid credentials'' & \textcolor{teal}{✓} \\
    3 & Compte rejeté (erreur 403) & Message ``Compte rejeté'' & \textcolor{teal}{✓} \\
    4 & Parsing erreur backend & Extraction correcte du message JSON & \textcolor{teal}{✓} \\
    \hline
  \end{tabularx}
\end{table}

\subsection{Guards de navigation -- 3 scénarios}

\begin{table}[H]
  \centering
  \caption{Scénarios de tests des guards RBAC}
  \label{tab:tests_guards}
  \rowcolors{2}{tablerow}{white}
  \begin{tabularx}{\textwidth}{|c|Y|Y|c|}
    \hline
    \thead{\#} & \thead{Action testée} & \thead{Comportement attendu} & \thead{Résultat} \\
    \hline
    1 & Accéder route -- utilisateur connecté & Accès autorisé & \textcolor{teal}{✓} \\
    2 & Accéder route -- non connecté & Redirection vers page de connexion & \textcolor{teal}{✓} \\
    3 & Accéder route MANAGER -- rôle DRIVER & Accès refusé & \textcolor{teal}{✓} \\
    \hline
  \end{tabularx}
\end{table}

\section{Tests de Performance -- k6}
\label{sec:ch8-performance}

\subsection{Objectif et framework}

Répondre à la question : \textit{``La plateforme LogiWay tient-elle si
50 chauffeurs et managers se connectent simultanément ?''}

\textbf{Framework :} k6 v2.2.0 (Grafana Labs -- open source).
Le login Keycloak est automatique au démarrage des tests.

\subsection{Concepts clés}

\begin{itemize}
  \item \textbf{VU (Virtual User)} : utilisateur simulé envoyant des requêtes
        en continu. 50~VUs = 50 utilisateurs simultanés.
  \item \textbf{p(95) = 38ms} : 95\% des utilisateurs reçoivent leur réponse
        en moins de 38ms. On utilise le percentile plutôt que la moyenne
        pour ignorer les pics extrêmes.
\end{itemize}

\subsection{Scénarios de test}

\begin{table}[H]
  \centering
  \caption{Scénarios de tests de performance}
  \label{tab:scenarios_perf}
  \rowcolors{2}{tablerow}{white}
  \begin{tabularx}{\textwidth}{|l|c|c|Y|}
    \hline
    \thead{Scénario} & \thead{Durée} & \thead{VUs max} & \thead{Objectif} \\
    \hline
    Smoke & 30s & 1 & Vérifier que l'API répond \\
    Load & 2min & 50 & Charge normale -- test principal \\
    Stress & 6min & 300 & Trouver le point de rupture \\
    Soak & 14min & 20 & Détecter les fuites mémoire \\
    \hline
  \end{tabularx}
\end{table}

\subsection{Fonctionnalités testées sous charge}

\begin{table}[H]
  \centering
  \caption{Fonctionnalités testées et seuils de performance}
  \label{tab:perf_seuils}
  \rowcolors{2}{tablerow}{white}
  \begin{tabularx}{\textwidth}{|Y|l|}
    \hline
    \thead{Fonctionnalité} & \thead{Seuil p(90)} \\
    \hline
    Liste des trajets de l'entreprise & < 1~500 ms \\
    Pauses réglementaires CE 561/2006 & < 5~000 ms \\
    Notifications temps réel SSE & < 1~500 ms \\
    Durée globale toutes requêtes p(95) & < 2~000 ms \\
    Taux d'erreur & < 5\% \\
    \hline
  \end{tabularx}
\end{table}

\subsection{Résultats mesurés -- Load Test (50 VUs, 2 minutes)}

\begin{table}[H]
  \centering
  \caption{Résultats du test de charge principal}
  \label{tab:resultats_perf}
  \rowcolors{2}{tablerow}{white}
  \begin{tabularx}{\textwidth}{|l|c|c|c|}
    \hline
    \thead{Métrique} & \thead{Seuil fixé} & \thead{Résultat mesuré} & \thead{Statut} \\
    \hline
    Durée globale p(95) & < 2~000 ms & \textbf{38 ms} & \textcolor{teal}{✓ PASSÉ} \\
    Trajets p(90) & < 1~500 ms & \textbf{30 ms} & \textcolor{teal}{✓ PASSÉ} \\
    Pauses CE 561/2006 p(90) & < 5~000 ms & \textbf{9 ms} & \textcolor{teal}{✓ PASSÉ} \\
    Notifications p(90) & < 1~500 ms & \textbf{39 ms} & \textcolor{teal}{✓ PASSÉ} \\
    Taux d'erreur & < 5\% & \textbf{0\%} & \textcolor{teal}{✓ PASSÉ} \\
    \hline
  \end{tabularx}
\end{table}

\begin{table}[H]
  \centering
  \caption{Indicateurs globaux du load test}
  \label{tab:indicateurs_perf}
  \rowcolors{2}{tablerow}{white}
  \begin{tabularx}{\textwidth}{|l|c|}
    \hline
    \thead{Indicateur} & \thead{Valeur} \\
    \hline
    Requêtes totales & $\sim$2~050 en 2 minutes \\
    Débit moyen & $\sim$17 requêtes/seconde \\
    Temps de réponse moyen & 22 ms \\
    Temps de réponse minimum & 4,5 ms \\
    Temps de réponse maximum & 246 ms \\
    VUs simultanés au pic & 50 \\
    Données reçues & 296 MB \\
    \hline
  \end{tabularx}
\end{table}

Le p(95) à \textbf{38~ms} est \textbf{52 fois inférieur} au seuil de
2~000~ms, confirmant que l'architecture Spring Boot + MySQL est correctement
dimensionnée pour un usage réel.

\section{Tests End-to-End -- Playwright}
\label{sec:ch8-e2e}

\subsection{Objectif et framework}

Simuler un vrai utilisateur humain naviguant dans l'application via
Chrome. Les tests E2E valident l'intégration complète :
Angular + Spring Boot + Keycloak + MySQL.

\begin{table}[H]
  \centering
  \caption{Frameworks des tests E2E}
  \label{tab:fw_e2e}
  \rowcolors{2}{tablerow}{white}
  \begin{tabularx}{\textwidth}{|l|Y|}
    \hline
    \thead{Outil} & \thead{Rôle} \\
    \hline
    Playwright 1.63.0 & Framework de contrôle de navigateur (Microsoft) \\
    Chromium 1243 & Navigateur Chrome utilisé pour l'exécution \\
    Global Setup & Login Keycloak unique partagé entre tous les tests \\
    \hline
  \end{tabularx}
\end{table}

\textbf{Architecture de session :} Le login Keycloak est effectué
\textbf{une seule fois} avant l'ensemble des tests. La session est
sauvegardée et injectée dans chaque test pour éviter les expirations.

\subsection{Groupe 1 : Authentification -- 4 tests}

\begin{table}[H]
  \centering
  \caption{Scénarios E2E du module authentification}
  \label{tab:e2e_auth}
  \rowcolors{2}{tablerow}{white}
  \begin{tabularx}{\textwidth}{|c|Y|Y|c|}
    \hline
    \thead{\#} & \thead{Action utilisateur} & \thead{Vérification} & \thead{Résultat} \\
    \hline
    1 & Afficher la page de connexion & Formulaire complet visible & \textcolor{teal}{✓} \\
    2 & Soumettre mauvais mot de passe & Message d'erreur affiché & \textcolor{teal}{✓} \\
    3 & Accéder sans session & Page de connexion accessible & \textcolor{teal}{✓} \\
    4 & Se connecter (MANAGER valide) & Redirection dashboard MANAGER & \textcolor{teal}{✓} \\
    \hline
  \end{tabularx}
\end{table}

\subsection{Groupe 2 : Navigation Dashboard -- 6 tests}

\begin{table}[H]
  \centering
  \caption{Scénarios E2E de navigation}
  \label{tab:e2e_nav}
  \rowcolors{2}{tablerow}{white}
  \begin{tabularx}{\textwidth}{|c|Y|Y|c|}
    \hline
    \thead{\#} & \thead{Page naviguée} & \thead{Vérification} & \thead{Résultat} \\
    \hline
    1 & Tableau de bord principal & Page chargée correctement & \textcolor{teal}{✓} \\
    2 & Gestion des Trajets & Pas de redirection vers signin & \textcolor{teal}{✓} \\
    3 & Gestion de la Flotte & Pas de redirection vers signin & \textcolor{teal}{✓} \\
    4 & Gestion des Réclamations & Pas de redirection vers signin & \textcolor{teal}{✓} \\
    5 & Profil utilisateur & Pas de redirection vers signin & \textcolor{teal}{✓} \\
    6 & Centre d'Alertes & Pas de redirection vers signin & \textcolor{teal}{✓} \\
    \hline
  \end{tabularx}
\end{table}

\subsection{Groupe 3 : Module Réclamations -- 5 tests}

\begin{table}[H]
  \centering
  \caption{Scénarios E2E du module réclamations}
  \label{tab:e2e_recla}
  \rowcolors{2}{tablerow}{white}
  \begin{tabularx}{\textwidth}{|c|Y|Y|c|}
    \hline
    \thead{\#} & \thead{Action utilisateur} & \thead{Vérification} & \thead{Résultat} \\
    \hline
    1 & Naviguer vers réclamations & Titre ``Gestion des Réclamations'' visible & \textcolor{teal}{✓} \\
    2 & Consulter les statistiques & 4 blocs : Total, Résolues, Rejetées, En cours & \textcolor{teal}{✓} \\
    3 & Consulter l'historique & Tableau avec données visible & \textcolor{teal}{✓} \\
    4 & Vérifier permissions MANAGER & Bouton ``Nouvelle réclamation'' visible & \textcolor{teal}{✓} \\
    5 & Ouvrir formulaire de création & Dialog Angular Material s'ouvre & \textcolor{teal}{✓} \\
    \hline
  \end{tabularx}
\end{table}

\section{Containerisation et déploiement Docker}
\label{sec:ch8-docker}

\subsection{Architecture Docker Compose}

La plateforme LogiWay est entièrement conteneurisée via Docker Compose
incluant les services suivants :

\begin{table}[H]
  \centering
  \caption{Services Docker Compose de LogiWay}
  \label{tab:docker_services}
  \rowcolors{2}{tablerow}{white}
  \begin{tabularx}{\textwidth}{|l|l|Y|}
    \hline
    \thead{Service} & \thead{Port} & \thead{Description} \\
    \hline
    mysql & 3306 & Base de données MySQL 8.0 \\
    keycloak & 8180 & Serveur d'authentification Keycloak 24 \\
    backend & 8080 & API Spring Boot 3.2.2 \\
    frontend & 4200/80 & Angular 19 servi par Nginx \\
    pause-ai & 5000 & Service Flask ML RandomForest \\
    reclamation-ai & 5001 & Service Flask IA système expert \\
    rag-service & 5003 & Service FastAPI RAG + Gemini \\
    \hline
  \end{tabularx}
\end{table}

\subsection{Démarrage de la plateforme complète}

La plateforme démarre en une seule commande :

\begin{tcolorbox}[ruleframe]
\texttt{docker-compose up -{}-build}
\end{tcolorbox}

Les variables de configuration sont centralisées dans le fichier
\texttt{.env} à la racine du projet pour une gestion sécurisée des
secrets et des paramètres d'environnement.

\section{Bilan global des tests}
\label{sec:ch8-bilan}

\begin{table}[H]
  \centering
  \caption{Bilan complet des tests LogiWay}
  \label{tab:bilan_tests}
  \rowcolors{2}{tablerow}{white}
  \begin{tabularx}{\textwidth}{|l|l|c|c|c|}
    \hline
    \thead{Couche} & \thead{Outil} & \thead{Tests} & \thead{Couverture} & \thead{Statut} \\
    \hline
    Backend Java & JaCoCo & 10 & 97,7\% lignes & \textcolor{teal}{✓} \\
    Python Toxicité & pytest + coverage.py & 7 & 99\% & \textcolor{teal}{✓} \\
    Python Sémantique & pytest + coverage.py & 5 & 99\% & \textcolor{teal}{✓} \\
    Python API & pytest-flask & 3 & 99\% & \textcolor{teal}{✓} \\
    Frontend Login & Jest + Istanbul & 4 & 100\% & \textcolor{teal}{✓} \\
    Frontend Guards & Jest + Istanbul & 3 & 100\% & \textcolor{teal}{✓} \\
    Performance Smoke & k6 & 1 & --- & \textcolor{teal}{✓} \\
    Performance Load & k6 & 1 & p(95)=38ms & \textcolor{teal}{✓} \\
    Performance Stress & k6 & 1 & --- & \textcolor{teal}{✓} \\
    Performance Soak & k6 & 1 & --- & \textcolor{teal}{✓} \\
    E2E Auth & Playwright & 4 & --- & \textcolor{teal}{✓} \\
    E2E Navigation & Playwright & 6 & --- & \textcolor{teal}{✓} \\
    E2E Réclamations & Playwright & 5 & --- & \textcolor{teal}{✓} \\
    \hline
    \textbf{Total} & & \textbf{60+} & & \textcolor{teal}{\textbf{✓ TOUS}} \\
    \hline
  \end{tabularx}
\end{table}

\textbf{Points forts de la stratégie de test :}
\begin{itemize}
  \item Couverture supérieure à 97\% sur toutes les couches de code
  \item 0\% de taux d'erreur sous 50 utilisateurs simultanés
  \item p(95) à 38~ms, 52 fois en dessous du seuil de 2~000~ms
  \item Tests E2E sans modification du code applicatif existant
  \item Login Keycloak automatique partagé entre tous les tests
  \item Règle CE~561/2006 testée sur 6 cas dont des tests paramétrés
\end{itemize}

\newpage


% ============================================================
% CHAPITRE 9 — RÉALISATIONS ET INTERFACES
% ============================================================
\chapitreavecplan{Réalisations et interfaces utilisateur}
\begin{plandechapitre}
  \planitem{Introduction}{sec:ch9-introduction}
  \planitem{Interface d'authentification}{sec:ch9-auth}
  \planitem{Interface SuperAdmin}{sec:ch9-superadmin}
  \planitem{Interface cartographique GPS temps réel}{sec:ch9-cartographie}
  \planitem{Interface Chauffeur -- Espace Cockpit}{sec:ch9-chauffeur}
  \planitem{Interface Chatbot et Génération de Rapports}{sec:ch9-chatbot}
  \planitem{Récapitulatif des interfaces réalisées}{sec:ch9-recap}
\end{plandechapitre}

\section{Introduction}
\label{sec:ch9-introduction}

Ce chapitre présente les interfaces réalisées pour chaque rôle,
illustrées par des captures d'écran et descriptions fonctionnelles.

\section{Interface d'authentification}
\label{sec:ch9-auth}

\begin{figure}[H]
  \centering
  \begin{minipage}[b]{0.48\textwidth}
    \centering
    \includegraphics[width=\textwidth]{interface_login.png}
    \caption*{Interface connexion -- thème sombre}
  \end{minipage}
  \hfill
  \begin{minipage}[b]{0.48\textwidth}
    \centering
    \includegraphics[width=\textwidth]{interface_login_light.png}
    \caption*{Interface connexion -- thème clair}
  \end{minipage}
  \caption{Interfaces de connexion et d'authentification de LogiWay}
  \label{fig:login}
\end{figure}

Fonctionnalités : login Keycloak, messages d'erreur explicites
(compte rejeté, désactivé, identifiants invalides), lien mot de passe oublié.

\section{Interface SuperAdmin}
\label{sec:ch9-superadmin}

\begin{figure}[H]
  \centering
  \begin{subfigure}[b]{0.48\textwidth}
    \centering
    \includegraphics[width=\textwidth]{dhaboard_superadmin.png}
    \caption{Vue d'ensemble KPIs}
    \label{fig:superadmin_dash1}
  \end{subfigure}
  \hfill
  \begin{subfigure}[b]{0.48\textwidth}
    \centering
    \includegraphics[width=\textwidth]{dhaboard_superadmin1.png}
    \caption{Graphiques analytiques}
    \label{fig:superadmin_dash2}
  \end{subfigure}
  \caption{Dashboard SUPERADMIN -- Vue globale multi-entreprises}
  \label{fig:superadmin_final}
\end{figure}

Navigation disponible : Tableau de Bord, Gérer les Comptes, Cartographie,
Analytique, Gestion Logistique, Centre d'Alertes, Messenger.

\section{Interface cartographique GPS temps réel}
\label{sec:ch9-cartographie}

\begin{figure}[H]
  \centering
  \includegraphics[width=0.9\textwidth]{Capture d'écran 2026-05-18 165130.png}
  \caption{Carte Leaflet -- suivi GPS temps réel, tracé OSRM, pauses IA}
  \label{fig:carto}
\end{figure}

Éléments affichés : tracé bleu (itinéraire OSRM), marqueur camion SVG
animé avec rotation selon le cap, marqueurs pauses colorés par type et
score IA, popup riche (vitesse, carburant, ETA, prochaine pause).

\section{Interface Chauffeur -- Espace Cockpit}
\label{sec:ch9-chauffeur}

\begin{figure}[H]
  \centering
  \begin{subfigure}[b]{0.48\textwidth}
    \centering
    \includegraphics[width=\textwidth]{dashboard_chauuffeur.png}
    \caption{Vue d'ensemble}
    \label{fig:chauffeur_dash1}
  \end{subfigure}
  \hfill
  \begin{subfigure}[b]{0.48\textwidth}
    \centering
    \includegraphics[width=\textwidth]{dashboard_chauuffeur1.png}
    \caption{Informations de mission}
    \label{fig:chauffeur_dash2}
  \end{subfigure}
  \caption{Espace Chauffeur -- Cockpit de mission}
  \label{fig:chauffeur_final}
\end{figure}

L'Espace Chauffeur affiche : entreprise de rattachement, manager superviseur,
véhicule assigné, mission actuelle (si EN\_COURS : destination, ETA,
bouton ``Voir sur la carte'').

\section{Interface Chatbot et Génération de Rapports}
\label{sec:ch9-chatbot}

\begin{figure}[H]
  \centering
  \includegraphics[width=0.85\textwidth]{chatbot.png}
  \caption{Interface chatbot RAG -- questions en français naturel}
  \label{fig:chatbot}
\end{figure}

L'interface affiche : conversation en bulles, historique persisté,
formulaire de génération de rapports avec sélection du format
(PDF/CSV/TXT) et téléchargement automatique.

\section{Récapitulatif des interfaces réalisées}
\label{sec:ch9-recap}

\begin{table}[H]
  \centering
  \caption{Récapitulatif des interfaces utilisateur de LogiWay}
  \label{tab:interfaces}
  \rowcolors{2}{tablerow}{white}
  \begin{tabularx}{\textwidth}{|l|l|Y|}
    \hline
    \thead{Interface} & \thead{Rôle} & \thead{Modules couverts} \\
    \hline
    Page de connexion & Tous & Login, reset MDP, vérification email \\
    Dashboard principal & SuperAdmin & KPIs globaux, graphiques, classements \\
    Gestion entreprises & SuperAdmin & CRUD, workflow validation \\
    Gestion secteurs & SuperAdmin & CRUD, assignation Manager \\
    Gestion utilisateurs & SuperAdmin & CRUD Manager/Chauffeur \\
    Console de Gestion & Manager & KPIs entreprise, graphiques hebdo \\
    Gérer la Flotte & Manager & Véhicules, affectation chauffeur \\
    Cartographie GPS & Manager/SuperAdmin & Carte Leaflet, OSRM, pauses IA \\
    Gestion Congés & Manager & Validation, rejet, historique \\
    Gestion Réclamations & Manager/SuperAdmin & Traitement, statuts \\
    Messenger & Tous & Discussions temps réel WebSocket \\
    Cockpit de Mission & Chauffeur & Mission, véhicule, secteur, Manager \\
    Mes Congés & Chauffeur & Demande, suivi, annulation \\
    Mes Réclamations & Chauffeur & Soumission avec validation IA \\
    Chatbot + Rapports & SuperAdmin/Manager & RAG, génération PDF/CSV/TXT \\
    Dashboard Pauses & Manager/SuperAdmin & KPIs sécurité, conformité CE 561/2006 \\
    \hline
  \end{tabularx}
\end{table}

\section{Conclusion}

Les interfaces de LogiWay ont été développées avec Angular~17 selon une
approche rôle-centrique. Le design en thème sombre, la réactivité SSE
et WebSocket, et les contrôles de sécurité backend garantissent une
expérience utilisateur fluide, sécurisée et adaptée à un usage professionnel.

\newpage

% ============================================================
% CHAPITRE 10 — CONCLUSION ET PERSPECTIVES
% ============================================================
\chapitreavecplan{Conclusion et perspectives}
\begin{plandechapitre}
  \planitem{Synthèse des réalisations}{sec:ch10-synthese}
  \planitem{Compétences acquises}{sec:ch10-competences}
  \planitem{Perspectives d'évolution -- Phase 2}{sec:ch10-perspectives}
  \planitem{Conclusion générale}{sec:ch10-conclusion}
\end{plandechapitre}

\section{Synthèse des réalisations}
\label{sec:ch10-synthese}

\begin{table}[H]
  \centering
  \caption{Bilan complet des modules réalisés}
  \label{tab:bilan_modules}
  \rowcolors{2}{tablerow}{white}
  \begin{tabularx}{\textwidth}{|l|l|Y|c|}
    \hline
    \thead{Sprint} & \thead{Module} & \thead{Fonctionnalités principales} & \thead{Avancement} \\
    \hline
    \multirow{4}{*}{Sprint 1} &
    Authentification & Keycloak OAuth2/JWT, reset MDP, RBAC & \textcolor{teal}{100\%} \\
    & Utilisateurs & CRUD Manager/Chauffeur, sync Keycloak & \textcolor{teal}{100\%} \\
    & Entreprises & Workflow EN\_ATTENTE$\rightarrow$ACTIF & \textcolor{teal}{100\%} \\
    & Secteurs + Véhicules & CRUD, affectation, isolation entreprise & \textcolor{teal}{100\%} \\
    \hline
    \multirow{2}{*}{Sprint 2} &
    Congés & Workflow EN\_ATTENTE/APPROUVE/REJETE & \textcolor{teal}{100\%} \\
    & Notifications SSE & Persistance MySQL + diffusion temps réel & \textcolor{teal}{100\%} \\
    \hline
    \multirow{2}{*}{Sprint 3} &
    Messagerie & Discussions WebSocket STOMP + historique & \textcolor{teal}{100\%} \\
    & Chatbot RAG & LangChain + Gemini + génération rapports & \textcolor{teal}{100\%} \\
    \hline
    \multirow{3}{*}{Sprint 4} &
    Trajets + GPS & Planification OSRM, suivi Multi-GNSS & \textcolor{teal}{100\%} \\
    & Pauses ML & RandomForest, CE 561/2006, affichage Leaflet & \textcolor{teal}{100\%} \\
    & Réclamations IA & Système expert toxicité + sémantique & \textcolor{teal}{100\%} \\
    \hline
    Sprint 5 &
    Dashboards KPI & SuperAdmin, Manager, Chauffeur, Pauses & \textcolor{teal}{100\%} \\
    \hline
    \multirow{2}{*}{Sprint 6} &
    Tests (60+) & JUnit 5, pytest, Jest, k6, Playwright & \textcolor{teal}{100\%} \\
    & Docker & Conteneurisation complète (7 services) & \textcolor{teal}{100\%} \\
    \hline
  \end{tabularx}
\end{table}

\section{Compétences acquises}
\label{sec:ch10-competences}

Ce projet m'a permis de maîtriser :

\begin{itemize}
  \item \textbf{Architecture sécurisée} : OAuth2/JWT Keycloak, Spring Boot
        Resource Server, cookies HttpOnly, protection XSS/CSRF.
  \item \textbf{Développement Full-Stack} : Spring Boot 3.x, Angular~17,
        JPA/Hibernate héritage JOINED, DTO pattern, pagination.
  \item \textbf{Communication temps réel} : WebSocket STOMP (GPS, messagerie),
        SSE (notifications), EventSource Angular.
  \item \textbf{Géolocalisation et cartographie} : Leaflet.js, OSRM poids lourds,
        Multi-GNSS GPS/Galileo/GLONASS.
  \item \textbf{Intelligence Artificielle} : système expert règles (Flask),
        RandomForest scikit-learn, RAG LangChain + Gemini, FastAPI.
  \item \textbf{Tests automatisés} : JUnit 5 + Mockito, pytest, Jest,
        k6 (performance), Playwright (E2E) -- 60+ tests, couverture > 97\%.
  \item \textbf{DevOps} : Docker, Docker Compose, conteneurisation
        multi-services, variables d'environnement.
  \item \textbf{Méthodologie Agile Scrum} : 6 sprints, backlog,
        revues, rétrospectives.
\end{itemize}

\section{Perspectives d'évolution -- Phase 2}
\label{sec:ch10-perspectives}

\subsection{Modèles Machine Learning avancés}

\begin{table}[H]
  \centering
  \caption{Modèles ML prévus pour la Phase 2}
  \label{tab:ml_phase2}
  \rowcolors{2}{tablerow}{white}
  \begin{tabularx}{\textwidth}{|l|l|Y|}
    \hline
    \thead{Module} & \thead{Algorithme} & \thead{Objectif} \\
    \hline
    Recommandation chauffeur & Régression logistique &
      Chauffeur optimal pour une mission donnée \\
    Prédiction fatigue avancée & Random Forest &
      Niveau de fatigue (0--100\%) selon historique \\
    Maintenance prédictive & XGBoost Regression &
      Probabilité de panne à 30 jours \\
    Anomalie carburant & Isolation Forest &
      Consommations anormales et cause probable \\
    \hline
  \end{tabularx}
\end{table}

\subsection{Application mobile}

Une application Android/iOS avec Angular/Ionic intégrera :
suivi GPS en arrière-plan, notifications push, messagerie mobile
et mode hors-ligne avec synchronisation.

\subsection{Tests d'intégration}

L'ajout de tests d'intégration avec Testcontainers (MySQL conteneurisé)
permettra de valider les interactions entre toutes les couches
(Controller + Service + Repository + MySQL réel) sans dépendance
à un environnement de déploiement.

\section{Conclusion générale}
\label{sec:ch10-conclusion}

Ce projet de développement chez \textbf{ExypnoTech Engineering Services}
a abouti à la conception d'une plateforme professionnelle de gestion de flotte
répondant aux enjeux opérationnels du transport routier tunisien.

La combinaison \textbf{Spring Boot~3.x + Angular~17 + Keycloak + MySQL},
enrichie de trois microservices d'intelligence artificielle
(système expert réclamations, RandomForest pauses, RAG chatbot),
produit une solution robuste, sécurisée et évolutive.

La stratégie de tests en 6 niveaux -- unitaires Backend (97,7\%),
Python IA (99\%), Frontend (100\%), performance k6 (38~ms p95, 0\% erreurs),
et E2E Playwright (15/15) -- garantit la qualité et la fiabilité
de l'ensemble de la plateforme.

Avec un avancement global de \textbf{100\%} des 6 sprints planifiés,
LogiWay démontre la faisabilité d'une plateforme intelligente de gestion
de flotte entièrement open-source, adaptée au contexte tunisien et
déployable à coût maîtrisé en gardant le contrôle des données.

\medskip
\begin{flushright}
  \textit{Kossay Boubaker,}\\
  \textit{Monastir, Septembre 2026}
\end{flushright}

\newpage

% ============================================================
% BIBLIOGRAPHIE
% ============================================================
\begin{thebibliography}{99}

\bibitem{springboot}
\textbf{Spring Boot 3.x Reference Documentation}\\
\url{https://docs.spring.io/spring-boot/docs/current/reference/html/}

\bibitem{springsec}
\textbf{Spring Security -- OAuth2 Resource Server}\\
\url{https://docs.spring.io/spring-security/reference/servlet/oauth2/resource-server/}

\bibitem{keycloak}
\textbf{Keycloak 24.x Documentation}\\
\url{https://www.keycloak.org/documentation}

\bibitem{angular}
\textbf{Angular 17 Documentation}\\
\url{https://angular.dev/overview}

\bibitem{mysql}
\textbf{MySQL 8.0 Reference Manual}\\
\url{https://dev.mysql.com/doc/refman/8.0/en/}

\bibitem{hibernate}
\textbf{Hibernate ORM 6 Documentation}\\
\url{https://hibernate.org/orm/documentation/}

\bibitem{osrm}
\textbf{OSRM -- Open Source Routing Machine}\\
\url{http://project-osrm.org/}

\bibitem{leaflet}
\textbf{Leaflet.js Documentation}\\
\url{https://leafletjs.com/reference.html}

\bibitem{sklearn}
\textbf{scikit-learn : RandomForestRegressor}\\
\url{https://scikit-learn.org/stable/modules/ensemble.html}

\bibitem{langchain}
\textbf{LangChain Documentation}\\
\url{https://python.langchain.com/docs/introduction/}

\bibitem{gemini}
\textbf{Google Gemini API Documentation}\\
\url{https://ai.google.dev/gemini-api/docs}

\bibitem{fastapi}
\textbf{FastAPI Documentation}\\
\url{https://fastapi.tiangolo.com/}

\bibitem{flask}
\textbf{Flask Documentation}\\
\url{https://flask.palletsprojects.com/}

\bibitem{reportlab}
\textbf{ReportLab Documentation}\\
\url{https://www.reportlab.com/docs/reportlab-userguide.pdf}

\bibitem{oauth2}
\textbf{OAuth 2.0 Authorization Framework -- RFC 6749}\\
\url{https://datatracker.ietf.org/doc/html/rfc6749}

\bibitem{ce561}
\textbf{Règlement (CE) No 561/2006 -- Temps de conduite et de repos}\\
\url{https://eur-lex.europa.eu/legal-content/FR/TXT/?uri=CELEX:32006R0561}

\bibitem{scrum}
\textbf{Scrum Guide 2020} -- K.~Schwaber, J.~Sutherland\\
\url{https://scrumguides.org/scrum-guide.html}

\bibitem{stomp}
\textbf{STOMP Protocol Specification 1.2}\\
\url{https://stomp.github.io/stomp-specification-1.2.html}

\bibitem{overpass}
\textbf{Overpass API Documentation}\\
\url{https://wiki.openstreetmap.org/wiki/Overpass_API}

\bibitem{junit5}
\textbf{JUnit 5 User Guide}\\
\url{https://junit.org/junit5/docs/current/user-guide/}

\bibitem{mockito}
\textbf{Mockito Documentation}\\
\url{https://javadoc.io/doc/org.mockito/mockito-core/latest/index.html}

\bibitem{jest}
\textbf{Jest Documentation}\\
\url{https://jestjs.io/docs/getting-started}

\bibitem{playwright}
\textbf{Playwright Documentation}\\
\url{https://playwright.dev/docs/intro}

\bibitem{k6}
\textbf{k6 Documentation -- Grafana Labs}\\
\url{https://k6.io/docs/}

\bibitem{jacoco}
\textbf{JaCoCo Java Code Coverage Library}\\
\url{https://www.jacoco.org/jacoco/trunk/doc/}

\bibitem{pytest}
\textbf{pytest Documentation}\\
\url{https://docs.pytest.org/en/stable/}

\end{thebibliography}

\end{document}
